% =====================================================================
%  A Partition-Theoretic Framework for Bioreactor Process Monitoring
%
%  Self-contained. Every concept used is defined within this document.
%  No companion manuscript is required.
% =====================================================================
\documentclass[11pt,twocolumn,a4paper]{article}

\usepackage[utf8]{inputenc}
\usepackage[T1]{fontenc}
\usepackage{lmodern}
\usepackage[a4paper,margin=2.5cm]{geometry}
\usepackage{amsmath,amssymb,amsthm,mathtools}
\usepackage{bm}
\usepackage{mathrsfs}
\usepackage{booktabs}
\usepackage{array}
\usepackage{enumitem}
\usepackage{microtype}
\usepackage{graphicx}
\usepackage{xcolor}
\usepackage{algorithm}
\usepackage{algpseudocode}
\usepackage[version=4]{mhchem}
\usepackage[round,sort&compress]{natbib}
\usepackage[colorlinks=true,linkcolor=blue!45!black,
            citecolor=blue!45!black,urlcolor=blue!45!black]{hyperref}
\usepackage{cleveref}

% ---- Theorem environments ----
\newtheoremstyle{thmstyle}{\topsep}{\topsep}{\itshape}{}{\bfseries}{.}{.5em}{}
\theoremstyle{thmstyle}
\newtheorem{theorem}{Theorem}[section]
\newtheorem{lemma}[theorem]{Lemma}
\newtheorem{proposition}[theorem]{Proposition}
\newtheorem{corollary}[theorem]{Corollary}
\theoremstyle{definition}
\newtheorem{definition}[theorem]{Definition}
\newtheorem{axiom}{Axiom}
\newtheorem{principle}[theorem]{Principle}
\newtheorem{example}[theorem]{Example}
\theoremstyle{remark}
\newtheorem{remark}[theorem]{Remark}
\newtheorem{assumption}[theorem]{Assumption}

\crefname{axiom}{Axiom}{Axioms}
\Crefname{axiom}{Axiom}{Axioms}

% ---- Shorthand ----
\newcommand{\R}{\mathbb{R}}
\newcommand{\N}{\mathbb{N}}
\newcommand{\Z}{\mathbb{Z}}
\newcommand{\kB}{k_{\mathrm{B}}}
\newcommand{\Depth}{\mathcal{M}}
\newcommand{\Sspace}{\mathcal{S}}
\newcommand{\Scoord}{\mathbf{s}}
\newcommand{\Sk}{s_{k}}
\newcommand{\St}{s_{t}}
\newcommand{\Se}{s_{e}}
\newcommand{\cut}{\partial}
\newcommand{\co}{\complement}
\newcommand{\dcat}{d_{\mathrm{cat}}}
\newcommand{\dC}{d_{\mathrm{C}}}
\newcommand{\floorc}{\beta}
\newcommand{\Gph}{\Gamma}
\newcommand{\V}{V}
\newcommand{\E}{E}
\newcommand{\wt}{w}
\newcommand{\med}{\mathsf{m}}
\newcommand{\Hcat}{H}
\newcommand{\taup}{\tau_{\mathrm{p}}}
\newcommand{\Rens}{R}
\newcommand{\abs}[1]{\lvert#1\rvert}
\newcommand{\norm}[1]{\lVert#1\rVert}
\newcommand{\eps}{\varepsilon}
\newcommand{\Reac}{\rho}

\title{\textbf{A Partition-Theoretic Framework for Bioreactor Process
Monitoring:\\[3pt] Cells as Finite Observers, Reactions as Localisable
Events, and Monitoring as Bounded Inference}}

\author{%
Kundai Farai Sachikonye\\
\small Technical University of Munich\\
\small Chair of Proteomics and Bioanalytics, School of Life Sciences\\
\small \texttt{kundai.sachikonye@wzw.tum.de}}

\date{\today}

\begin{document}
\maketitle

% =====================================================================
\begin{abstract}
\noindent
Industrial and laboratory bioreactors are monitored today through a
patchwork of sensors---dissolved oxygen, pH, off-gas, spectroscopy,
capacitance---whose readings are fused by empirical soft sensors and
multivariate statistics. This works, but it rests on no unified account
of \emph{what a bioreactor measurement is}, \emph{how precise it can be
in principle}, or \emph{when the monitoring system is entitled to stop
and report}. We develop such an account from a single premise: that a
bioreactor is a finite population of cells, each of which is a
\emph{finite observer}---a bounded system that can distinguish only
finitely many states at any instant. From this premise, treated with the
mathematics of finite weighted graphs, information theory, and
measure-preserving dynamics, we derive a complete monitoring framework in
which four facts appear as theorems rather than assumptions. First, every
individuation of a process state carries a strictly positive
\emph{resolution floor}: no monitor returns an exact value, only a bounded
region, and this is a property of the finite medium, not a defect of the
instrument. Second, the state of the reaction network can be reconstructed
from \emph{partial} observation by a provably convergent
\emph{trajectory-completion} procedure that propagates measured quantities
through the network's own conservation and thermodynamic constraints.
Third, an individual biochemical reaction is not merely ``occurring
somewhere'' but is a localisable event, whose position and time are
recoverable to nanometre-scale precision by intersecting the
characteristic arrival-time surfaces of several physically distinct
propagation modalities that every reaction excites simultaneously. Fourth,
because a state is individuated by contrast against the whole rather than by
finding a cell that instantiates it, the framework can monitor
\emph{coincidence nodes}---processes that occur only under a conjunction (or
an ordered sequence) of conditions that no single cell need ever satisfy at
once---reconstructing states that are true of the population as a structured
whole while being present, at any instant, in no member of it; such
processes are provably invisible to per-cell and bulk-average monitoring.
A swarm of monitoring agents phase-locked to a common cellular clock
coordinates at zero communication overhead, and its loss of coherence is
itself a first-class fault signal. We give an explicit, honest stopping
rule---\emph{closure}---that is strictly stronger than crossing a
confidence threshold and that terminates every monitoring query either in
a single reported state or in an explicit declaration that the evidence
is contested. We establish, as a structural result, that process
monitoring and process control must be kept as separate frameworks: a
system that acts on its own observations cannot serve as an unbiased
observer of the process it acts on, and the monitoring floor forbids the
point-precision that a fused controller would presume. The framework's
worst-case output is therefore never a wrong action but detailed,
uncertainty-bounded information about the process. We specify the complete
monitoring pipeline as an algorithm, state its complexity and correctness
guarantees, and describe how each abstract quantity maps to a physical
bioreactor observable. Every result is proved from the stated premises and
additionally verified on explicitly constructed instances, where a suite of
seven experiments comprising thirty-four checks finds every measured
quantity in agreement with the predicted bound; the paper is self-contained.
\end{abstract}

\noindent\textbf{Keywords:} bioprocess monitoring; process analytical
technology; finite observer; resolution floor; partition depth;
categorical distance; trajectory completion; multimodal localisation;
oscillator synchronisation; closure; monitor--control separation.

\tableofcontents

% =====================================================================
\section{Introduction}
\label{sec:intro}
% =====================================================================

\subsection{The monitoring problem}

A bioreactor is a vessel in which a population of living cells converts
substrate into product under controlled physical conditions. Whether the
product is a therapeutic protein, a small-molecule metabolite, biomass, or
a gas, the economic and regulatory value of the process depends on knowing,
at each instant, \emph{what the cells are doing}. This is the monitoring
problem, and it is distinct from the control problem of \emph{making} the
cells do something: monitoring produces knowledge, control produces action.
The distinction is not merely organisational; we shall argue
(\Cref{sec:separation}) that it is forced by the structure of the problem
itself.

Contemporary practice equips the vessel with a heterogeneous sensor
array---dissolved-oxygen and pH probes, temperature and pressure
transducers, off-gas mass spectrometry, dielectric spectroscopy, optical
density and fluorescence, and increasingly Raman and near-infrared
spectroscopy \citep{lourenco2012bioreactor,sonnleitner2013automated}---and
fuses their outputs through \emph{soft sensors}: empirical or
semi-mechanistic models that infer unmeasured quantities (biomass, specific
growth rate, metabolite concentrations) from the measured ones
\citep{luttmann2012soft,mandenius2015bioprocess}. This activity is
codified, for pharmaceutical manufacturing, under the banner of Process
Analytical Technology \citep{fda2004pat}. The statistical backbone is
multivariate process monitoring---principal component analysis, partial
least squares, and their fault-detection extensions
\citep{wold1987principal,qin2012survey,venkatasubramanian2003review}.

This machinery is effective and mature. What it lacks is a first-principles
answer to three questions that become pressing precisely when monitoring
must be trustworthy:

\begin{enumerate}[leftmargin=1.6em,itemsep=3pt]
\item \textbf{What is a bioreactor measurement, exactly?} A dissolved-oxygen
reading is a number; but the quantity of interest---``the metabolic state
of the culture''---is not a number the probe reports. What is the
mathematical object that a measurement individuates, and what does it mean
for two differently-instrumented readings to report the same underlying
state?
\item \textbf{How precise can monitoring be, in principle?} Every soft
sensor carries an error bar, but these are estimated empirically. Is there
a floor---a precision that no instrument, however ideal, can beat---set by
the physics of the culture rather than by the quality of the hardware?
\item \textbf{When is the monitoring system entitled to stop and report?}
A monitor that fuses many sensors must decide when it has seen enough. The
usual answer, ``when a confidence score crosses a threshold,'' is
demonstrably too weak: a single self-consistent sensor can report high
confidence while a second, independent sensor would have reported something
entirely different. What is the correct stopping rule?
\end{enumerate}

\subsection{The premise and the approach}

We answer all three questions from a single premise. A bioreactor is a
finite population of cells; a cell is a bounded physical system; and a
bounded physical system, observed at finite resolution, can be in only
finitely many distinguishable states at any instant. We call such a system
a \emph{finite observer}. The premise is not a model choice but an
observation: cells have finite volume, finite conformational freedom,
finite copy numbers of each molecular species, and they operate against a
substrate reservoir and an environment they can never fully resolve.

From this premise we build, using only finite weighted graphs, information
theory, and measure-preserving dynamics, a monitoring framework whose
central facts are theorems:

\begin{itemize}[leftmargin=1.6em,itemsep=3pt]
\item \textbf{A resolution floor} (\Cref{sec:floor}). Because a cell can
never complete the comparison of its own state against an inexhaustible
environment, every state it individuates carries a strictly positive
irreducible residual. Monitoring inherits this floor: no monitor returns a
point, only a bounded region. The floor is derived, not posited, and it is
a property of the finite culture, not of the sensor.
\item \textbf{Trajectory completion} (\Cref{sec:completion}). The full
state of the metabolic reaction network can be reconstructed from partial
sensor data by iterating three constraint classes---mass balance,
thermodynamic consistency, and kinetic feasibility---to a unique fixed
point, guaranteed to exist and to be reached by a contraction argument.
\item \textbf{Reaction localisation} (\Cref{sec:localisation}). An
individual biochemical reaction excites disturbances in several physically
distinct propagation modalities at once. Because these modalities travel at
different speeds and obey different laws, the intersection of their
arrival-time surfaces pins the reaction's location and time; adding
modalities multiplies resolving power.
\item \textbf{Coincidence nodes} (\Cref{sec:coincidence}). Because a state
is fixed by contrast against the whole, not by finding a cell that
instantiates it, the framework monitors processes that occur only under a
conjunction---or an ordered sequence---of conditions no single cell need
ever satisfy at once. It thereby reconstructs states true of the population
as a structured whole yet present in no member of it, which per-cell and
bulk-average methods provably cannot detect.
\item \textbf{Swarm coordination} (\Cref{sec:swarm}). A distributed array
of monitoring agents that phase-lock to a common cellular oscillation
coordinates at zero communication overhead once coherent, and the loss of
that coherence is itself a first-class, early fault signal.
\item \textbf{Closure} (\Cref{sec:closure}). A monitoring query terminates
correctly when no further available evidence can change its outcome---a
criterion strictly stronger than a confidence threshold---and reports
either a single state or an explicit \emph{decline}: the finding that the
evidence is contested.
\end{itemize}

We then show (\Cref{sec:separation}) that these results jointly force the
separation of monitoring from control, and we assemble the whole into an
explicit pipeline (\Cref{sec:pipeline}) with stated complexity and
correctness. \Cref{sec:mapping} maps every abstract quantity to a physical
bioreactor observable, and \Cref{sec:discussion} states scope and
limitations.

\subsection{What is and is not claimed}

We claim a coherent, self-contained mathematical framework and a monitoring
algorithm that follows from it, together with an account of the precision
limits and stopping conditions of bioreactor monitoring. We do \emph{not}
claim a control law, a completed software system, or experimental
validation on a physical bioreactor; those are separate undertakings, and
the separation of monitoring from control is itself one of our results.
Where a term carries an evocative reading---``observer,'' ``reaction,''
``catalyst''---it names a specific mathematical object defined in the text,
and no theorem depends on any reading beyond that definition. Interpretive
remarks are marked as such and are never load-bearing.

\subsection{Organisation}

\Cref{sec:observer} fixes the finite-observer premise and the contact-graph
formalism. \Cref{sec:floor} derives the resolution floor.
\Cref{sec:sspace} introduces the S-entropy state space and partition depth,
the coordinates in which the culture's state is represented.
\Cref{sec:reactions} models reactions as partition operations and derives
their rates. \Cref{sec:completion} gives the trajectory-completion
algorithm and its convergence proof. \Cref{sec:localisation} develops
multimodal reaction localisation. \Cref{sec:swarm} treats swarm
coordination. \Cref{sec:coincidence} establishes the monitoring of
coincidence nodes---states true of the population but present in no cell.
\Cref{sec:closure} defines closure and the two ways a query ends. \Cref{sec:separation} proves the monitor--control separation.
\Cref{sec:pipeline} assembles the complete pipeline. \Cref{sec:mapping}
gives the physical mapping. \Cref{sec:validation} reports the numerical
validation suite. \Cref{sec:discussion} concludes.

% =====================================================================
\section{The Cell as a Finite Observer}
\label{sec:observer}
% =====================================================================

We fix the premise of the paper and the mathematical objects it produces.
Nothing in this section is specific to biology beyond the observation that
cells satisfy the premise.

\subsection{Bounded systems and finite distinguishability}

\begin{axiom}[Finite observer]
\label{ax:observer}
A \emph{finite observer} is a physical system that (i) occupies a bounded
region of its state space and (ii) at any instant can distinguish only
finitely many states of that region. A cell in a bioreactor is a finite
observer: its molecular configuration lies in a bounded region (finite
volume, finite energy, finite copy numbers), and its capacity to
distinguish configurations---of itself and of its surroundings---is limited
by finite molecular resources.
\end{axiom}

The two clauses are physically independent and both are needed.
Boundedness alone permits a continuum of indistinguishable states; finite
distinguishability alone permits an unbounded catalogue. Together they say:
a finite observer partitions its bounded world into finitely many
\emph{cells of distinguishability}, and everything it can know is which cell
it, or another part of the world, occupies.

\begin{definition}[The medium]
\label{def:medium}
The \emph{medium} of a finite observer is the totality of what it is
individuated against: the reactor contents, the environment, and the
substrate reservoir, taken as one undifferentiated whole prior to any
distinction. We denote the medium $\med$. The medium is
\emph{non-completable}: for any finite collection of states the observer
has already distinguished, there remains a further state of the medium not
yet distinguished. This is an order condition---no last stage of
distinction---not a claim about the cardinality of the medium.
\end{definition}

Non-completability is the formal content of the everyday fact that a cell
in a reactor is never in possession of a complete inventory of its
surroundings: there is always another molecule, another gradient, another
fluctuation not yet resolved. We make no stronger assumption.

\subsection{Contact graphs}

We represent what a finite observer has distinguished as a finite weighted
graph, following the standard theory of network cuts and flows
\citep{diestel2017graph,fordfulkerson1956maximal}.

\begin{definition}[Finite weighted graph]
\label{def:graph}
A \emph{finite weighted graph} is a triple $\Gph=(\V,\E,\wt)$ with $\V$ a
finite non-empty vertex set, $\E\subseteq\binom{\V}{2}$ a set of undirected
edges, and $\wt\colon\E\to\R_{>0}$ a strictly positive weight. We take
$\Gph$ connected unless stated, and simple (no loops or multi-edges).
\end{definition}

\begin{definition}[Contact graph]
\label{def:contact}
A \emph{contact graph} is a finite weighted graph
$\Gph=(\V,\E,\wt)$ together with a distinguished vertex $\med\in\V$, the
\emph{medium}, adjacent to every other vertex. Vertices $v\neq\med$ are
\emph{parts}: individuated units---a molecular species, a metabolite pool,
a reaction intermediate, a measured quantity. An edge $\{u,v\}$ between two
parts is a \emph{contact}, and its weight $\wt(\{u,v\})$ is the
\emph{cost of the distinction} separating $u$ from $v$: how much must be
resolved to tell them apart.
\end{definition}

The single medium vertex encodes a deliberate structural commitment: the
theory attaches no privileged status to where information physically
resides. A quantity read from an in-line probe and a quantity inferred from
off-gas analysis are, prior to individuation, both contents of $\med$; what
distinguishes a part from the medium is the \emph{cost} of individuating
it (\Cref{def:sepcost}), never its physical location. This is what will
license, later, treating chemically, optically, and electrically sensed
observations as uniform inputs to one inference (\Cref{sec:localisation}).

\begin{definition}[Cut and separation cost]
\label{def:sepcost}
For $S\subseteq\V$, the \emph{cut} of $S$ is
$\cut(S)=\{\,\{u,w\}\in\E : u\in S,\ w\notin S\,\}$, and its weight is
$\wt(\cut(S))=\sum_{e\in\cut(S)}\wt(e)$. The \emph{separation cost} of a
part $v$ is the minimum weight of a cut placing $v$ on one side and the
medium on the other,
\[
\sigma(v)\;=\;\min_{\substack{S\ni v\\ \med\notin S}}\wt(\cut(S)).
\]
Because $\V$ is finite the minimum is attained, and it is computable
exactly by maximum flow \citep{fordfulkerson1956maximal}.
\end{definition}

The separation cost is the graph-theoretic content of ``how hard it is to
tell $v$ apart from everything else.'' A metabolite fixed by a single
uncontested probe reading has low separation cost; a state that can be
established only by triangulating many weakly corroborating signals has
high separation cost. This quantity is the backbone of everything below.

\subsection{Individuation is by negation}

A finite observer with no privileged, pre-marked answer must fix a part by
what it is \emph{not}. We record this because it dictates the logical form
of every measurement in the framework.

\begin{definition}[No privileged part]
\label{def:noprivilege}
A determining rule has \emph{no privileged part} if the data available to
it are invariant under the weighted automorphisms of $\Gph$, so that no
part is singled out by the data alone.
\end{definition}

\begin{theorem}[Individuation by negation]
\label{thm:negation}
Let $\Gph$ have no privileged part. Then any rule that fixes a set of parts
$U\subseteq\V$ without naming a member of $U$ in advance fixes it as the
complement of its negation set, $U=\V\setminus\co U$ with $\co U:=\V\setminus
U$ specified first; and conversely every admissible negation set determines
a unique $U$. No rule that names a member of $U$ escapes an infinite
regress of such namings.
\end{theorem}

\begin{proof}
\emph{Necessity.} To fix $U$ is to decide, for each part, whether it lies
in $U$ or in the remainder. A rule that fixes $U$ by asserting membership
of a distinguished part must name that part; since the data are
automorphism-invariant (\Cref{def:noprivilege}), no part is named by the
data, so the rule must itself supply the name. That name must in turn be
fixed---given in advance (excluded) or by a prior rule facing the same
dichotomy. The regress halts only at a rule naming no member of $U$:
specifying $\co U=\V\setminus U$ and taking $U$ as the remainder, which
names only ``the rest,'' a datum the medium already supplies.
\emph{Sufficiency.} Complementation $U\mapsto\V\setminus U$ is an
involution on the finite set $\V$, so a given $\co U\subsetneq\V$ with
$\V\setminus\co U\neq\varnothing$ determines $U$ uniquely.
\end{proof}

\begin{corollary}[A part is fixed by its non-neighbours]
\label{cor:nonneighbours}
A single part $v$ is determined without a selector by the partition of the
rest of the medium into what it contacts and what it does not: its identity
is carried by this negation structure, not by any intrinsic label.
\end{corollary}

\begin{remark}[The monitoring reading]
\label{rem:monitoring-negation}
\Cref{thm:negation} is why a bioreactor measurement is properly a
\emph{contrast}, not an absolute. A dissolved-oxygen probe does not report
``oxygen''; it reports oxygen individuated against everything oxygen is not
in that medium. Two probes of different physical principle that individuate
the same part against the same medium report the same state, even when
their raw signals share no units. This is the sense in which the framework
treats heterogeneous sensors as commensurable: they are commensurable when,
and only when, they individuate the same part of the same medium. Nothing
below depends on any reading of ``negation'' beyond the graph statement.
\end{remark}

% =====================================================================
\section{The Resolution Floor of Monitoring}
\label{sec:floor}
% =====================================================================

We now derive the central quantitative fact of the framework: every
monitored quantity carries a strictly positive irreducible uncertainty,
forced by the finiteness of the observer and the non-completability of the
medium. It is not posited as an error budget; it is a theorem.

\subsection{Identification and its residual}

\begin{definition}[Identification residual]
\label{def:residual}
To \emph{identify} a part $v$ is to determine it by its minimum cut against
the medium (\Cref{def:sepcost}). The identification is \emph{complete} if
every part of the whole medium has been brought to bear on that cut. The
\emph{identification residual} $\varrho(v)\ge0$ is the boundary content a
complete identification would still remove that a finite-stage
identification cannot; $\varrho(v)=0$ iff the identification is complete.
\end{definition}

\begin{theorem}[Floor from non-completability]
\label{thm:floor-derived}
Let the medium be non-completable (\Cref{def:medium}) and let $v$ be a part
($v\neq\med$, and $v$ does not exhaust $\V$). Then the identification of
$v$ cannot be completed at any finite stage, and $\varrho(v)>0$.
\end{theorem}

\begin{proof}
Suppose $\varrho(v)=0$, so the identification of $v$ is complete at some
finite stage. By \Cref{def:residual}, completeness means every part of the
medium has entered the cut of $v$ against $\med$; equivalently, the finite
stage already contains the whole medium. This contradicts
non-completability, which guarantees a further part not yet present.
Hence no finite stage completes the identification, and $\varrho(v)>0$.
\end{proof}

\begin{theorem}[The resolution floor]
\label{thm:floor}
There is a constant $\floorc>0$, the \emph{resolution floor}, such that
$\wt(e)\ge\floorc$ for every contact $e$ and $\sigma(v)\ge\floorc$ for
every part $v$. No two parts are separated at zero cost, and no part is
individuated for free.
\end{theorem}

\begin{proof}
Set $\floorc:=\inf_v\varrho(v)$. By \Cref{thm:floor-derived},
$\varrho(v)>0$ for every part; the residual is the irreducible boundary
content of the identification of $v$ and is carried by the edges of its
minimum cut, so each such edge has weight at least the per-part residual,
hence at least $\floorc$. Every contact edge $\{u,v\}$ is a boundary edge
of the identification of $u$ (it lies in $\cut(\{u\})$), so $\wt(e)\ge
\floorc$ for all $e$. Connectivity forces $\cut(S)\neq\varnothing$ for
every proper nonempty $S$, so $\sigma(v)\ge\floorc>0$.
\end{proof}

\begin{corollary}[No point measurement]
\label{cor:no-point}
No monitor returns a part separated from the medium by a cut of weight
zero. Every monitored quantity resolves to a region of the state space
whose width is bounded below by the floor $\floorc$; a monitor that appears
to answer ``exactly'' has, in fact, resolved to such a region and no
smaller.
\end{corollary}

\begin{proof}
Immediate from \Cref{thm:floor}: a zero-width answer would require
$\sigma(v)=0<\floorc$.
\end{proof}

\subsection{The floor is a property of the culture, not the sensor}

\begin{remark}[Why the floor is not instrument error]
\label{rem:floor-not-error}
\Cref{cor:no-point} is often mistaken for a statement about sensor
imperfection. It is not. Even an ideal instrument cannot deliver a
zero-residual reading, because delivering one would require completing a
comparison against a medium that---by the non-completability of a finite
culture in an inexhaustible environment---can never be completed. The floor
is the trace the non-completable medium leaves on every finite measurement.
Improving the hardware lowers the practical error toward the floor; it
cannot cross it. This reframes the reporting of monitoring results: an
honest monitor reports a region of width $\ge\floorc$ and treats that width
as physics, not as an apology for imperfect sensing.
\end{remark}

\begin{proposition}[Refinement lowers but never abolishes the floor]
\label{prop:refine}
Let $\Gph'$ be obtained from $\Gph$ by adding parts and contacts (a finer
description obtained from more or better sensors), each new contact of
weight $\ge\floorc$. Then the realised floor
$\floorc_{\Gph'}:=\min_{v}\sigma_{\Gph'}(v)$ satisfies
$\floorc\le\floorc_{\Gph'}$ up to the added structure, remains strictly
positive at every finite stage, and reaches zero only in the
non-realisable limit of an exhausted medium.
\end{proposition}

\begin{proof}
$\Gph'$ is again a finite contact graph with all weights $\ge\floorc$, so
\Cref{thm:floor} applies and $\sigma_{\Gph'}(v)\ge\floorc>0$ for every part.
Adding structure can only introduce parts with smaller separation cost, but
never below $\floorc$. The floor vanishes only if the medium is exhausted,
which non-completability forbids at every finite stage.
\end{proof}

\begin{remark}[Consequence for reporting]
\label{rem:report}
\Cref{prop:refine} gives the monitoring engineer an exact expectation:
adding sensors sharpens the reported region and can only help, but there is
a diminishing return and a hard floor. A monitor that claims sub-floor
precision is reporting an artefact. \Cref{sec:closure} converts this into
a stopping rule: keep adding evidence until further evidence cannot change
the reported region---which, by the floor, is a well-defined and reachable
state.
\end{remark}

\begin{figure*}[!htbp]
    \centering
    \includegraphics[width=\textwidth]{figures/panel1_floor.png}
    \caption{\textbf{The Resolution Floor: No Individuation Is Free.}
    (\textbf{A})~Separation cost $\sigma(v)$ on a log $y$-axis
    ($10^{0}$ to ${\sim}10^{2}$) versus part count on the linear $x$-axis
    ($3$ to $14$), one point per part over $\sim\!2{,}600$ parts drawn from
    $300$ random contact graphs (steel-blue), with the floor $\floorc=1$
    marked (coral dashed). Every measured separation lies on or above the
    floor---the smallest observed ratio $\sigma(v)/\floorc$ is $1.01$ and the
    mean is $26.1$---the empirical content of \Cref{thm:floor}: no part is
    individuated from the medium at zero cost, so no monitored quantity
    resolves to a point (\Cref{cor:no-point}).
    (\textbf{B})~Distribution of the dimensionless ratio $\sigma(v)/\floorc$
    over all parts (sea-green), with the threshold $\sigma/\floorc=1$ marked
    (coral dashed). The histogram is supported entirely at and above $1$:
    across every constructed instance not one separation falls below the
    floor, the no-sharp-cut clause of \Cref{thm:floor} with no exceptions.
    (\textbf{C})~Minimum separation ratio $\min_v\sigma(v)/\floorc$ on the
    linear $y$-axis versus medium size on a log-$2$ $x$-axis
    ($2^{2}$ to $2^{6}$), floor marked (coral dashed). The realised floor
    remains strictly above $1$ as the medium grows from four to sixty-four
    parts and rises with size, confirming \Cref{prop:refine}: enlarging the
    description narrows nothing below the floor and never abolishes it.
    (\textbf{D})~Three-dimensional surface of the per-instance minimum
    separation over part count ($3$ to $12$) and floor $\floorc$
    ($1$ to $6$), viridis colourmap, viewing angle elevation $22^{\circ}$
    azimuth $-58^{\circ}$. The surface is strictly positive everywhere and
    rises with $\floorc$: the floor is a receiver-intrinsic property that
    tracks the coarsest admissible distinction, as \Cref{thm:floor-derived}
    and \Cref{thm:floor} require.}
    \label{fig:panel1}
\end{figure*}

% =====================================================================
\section{State Coordinates: S-Entropy Space and Partition Depth}
\label{sec:sspace}
% =====================================================================

The contact graph records what has been distinguished; we now fix the
coordinates in which a monitored state is reported and the scalar that
measures how much structure a state carries. These are the working
variables of the monitoring pipeline.

\subsection{The S-entropy state space}

A finite observer distinguishes states along three independent axes of
uncertainty. We normalise each to $[0,1]$ and take their product.

\begin{definition}[S-entropy coordinates]
\label{def:sspace}
The \emph{S-entropy state space} is the unit cube
$\Sspace=[0,1]^3$. A state is a point $\Scoord=(\Sk,\St,\Se)$ where
\begin{itemize}[leftmargin=1.4em,nosep]
\item $\Sk$ (\emph{knowledge entropy}) measures uncertainty in \emph{which}
state the system occupies---identification uncertainty;
\item $\St$ (\emph{temporal entropy}) measures uncertainty in the
\emph{timing} of the system's transitions;
\item $\Se$ (\emph{evolution entropy}) measures uncertainty in the
system's \emph{trajectory}---where it is heading.
\end{itemize}
Each coordinate is a normalised information measure
\citep{shannon1948mathematical,jaynes1957information}; for a quantity $x$
with reference scale $x_0$ and fluctuation range, the normalisation maps
observed spread to $[0,1]$ monotonically.
\end{definition}

\begin{proposition}[$\Sspace$ is a bounded metric space with an intrinsic
metric]
\label{prop:metric}
$\Sspace=[0,1]^3$ is compact and connected, and carries the Euclidean
metric
\[
d(\Scoord_1,\Scoord_2)=\sqrt{(\Sk_1-\Sk_2)^2+(\St_1-\St_2)^2+(\Se_1-\Se_2)^2},
\]
computable in constant time from the coordinates of two states, independent
of how many states the system has visited.
\end{proposition}

\begin{proof}
Compactness and connectedness of $[0,1]^3$ are standard; the Euclidean
distance is a metric computable from two triples of reals in $O(1)$
arithmetic operations, with no dependence on the size of any state
catalogue.
\end{proof}

\begin{remark}[Why an intrinsic metric matters operationally]
\label{rem:intrinsic}
\Cref{prop:metric} means the monitor never needs a precomputed table of
pairwise state distances: the distance between any two culture states is
read off their coordinates directly. For a monitoring system tracking a
culture through a large state space, this is the difference between storing
a distance matrix (quadratic in the number of states) and storing only the
coordinates (linear). Distances---and hence the search for the state
consistent with the sensors---are computed on demand.
\end{remark}

\subsection{Partition depth}

We need a scalar that says how much distinguishing structure a state
carries. This is the quantity whose gradients, we shall see, drive reaction
dynamics.

\begin{definition}[Partition depth]
\label{def:depth}
The \emph{partition depth} $\Depth$ of a state is the amount of
hierarchical distinction required to specify it uniquely among the
alternatives:
\[
\Depth\;=\;\sum_{i}\log_b(k_i),
\]
where $k_i$ is the branching factor at level $i$ of the specification and
$b$ is the encoding base. A coarse distinction has low depth; a fine one
has high depth.
\end{definition}

\begin{proposition}[Depth is entropy]
\label{prop:depth-entropy}
The thermodynamic entropy associated with a state of partition depth
$\Depth$ is $S_{\mathrm{p}}=\kB\Depth\ln b$, so
$\Delta S_{\mathrm{p}}=\kB\ln b\cdot\Delta\Depth$.
\end{proposition}

\begin{proof}
Partition depth counts the $b$-ary distinctions needed to specify a state;
each distinction is $\ln b$ nats of information. By the Boltzmann relation
$S=\kB\ln W$, a state distinguished among $W=b^{\Depth}$ alternatives
carries entropy $\kB\ln(b^{\Depth})=\kB\Depth\ln b$
\citep{shannon1948mathematical,cover2006elements}.
\end{proof}

\begin{corollary}[Depth is bounded below and above]
\label{cor:depth-bounds}
For any finite observer, $\log_b 2\le\Depth\le\log_b(\Omega/h_0)$, where the
lower bound is one distinction (existence as a distinct state) and the
upper bound is the total distinguishable-state count $\Omega/h_0$ of the
bounded region at resolution $h_0$. A state of depth $0$ is unobservable.
\end{corollary}

\begin{proof}
Observation requires at least one distinction, giving the lower bound. A
bounded region resolved at scale $h_0$ has at most $\Omega/h_0$
distinguishable states, and enumerating them needs depth
$\log_b(\Omega/h_0)$, giving the upper bound
\citep{cover2006elements}.
\end{proof}

\subsection{Categorical distance}

Finally, we need the distance a process must travel through the space of
distinctions to get from one state to another. This is the quantity that
will set reaction rates and, later, resolving power.

\begin{definition}[Categorical distance]
\label{def:catdist}
The \emph{categorical distance} between two states is the number of
partition boundaries that must be crossed to pass from one to the other:
\[
\dcat(A,B)\;=\;\text{number of partition boundaries between } A \text{ and } B.
\]
Equivalently, if states are labelled by discrete partition coordinates,
$\dcat$ is the $\ell^1$ distance between their coordinate labels.
\end{definition}

\begin{proposition}[$\dcat$ is a metric]
\label{prop:catdist-metric}
$\dcat$ is non-negative, symmetric, zero exactly on identical states, and
satisfies the triangle inequality.
\end{proposition}

\begin{proof}
All four properties are inherited from the $\ell^1$ distance on the
discrete coordinate labels; non-negativity, symmetry, and identity are
immediate, and the triangle inequality is that of $\ell^1$.
\end{proof}

\begin{remark}[Two distances, two roles]
\label{rem:two-distances}
The framework uses two distances, and it is worth keeping them apart. The
Euclidean $d$ on $\Sspace$ (\Cref{prop:metric}) measures how far apart two
\emph{culture states} are for the purpose of searching and navigation. The
categorical $\dcat$ (\Cref{def:catdist}) measures how many \emph{distinct
steps} a reaction must take, and it is this that sets reaction rates
(\Cref{sec:reactions}) and, crucially, is independent of physical
distance---a fact we exploit for opacity-independent measurement
(\Cref{sec:localisation}).
\end{remark}

% =====================================================================
\section{Reactions as Localisable Partition Events}
\label{sec:reactions}
% =====================================================================

We now model the objects the monitor is ultimately watching: the
biochemical reactions of the culture. We derive their rates from the
partition structure, cast the reaction network as a circuit whose node
potentials are information-theoretic, and establish the fact that motivates
external monitoring in the first place---that a cell cannot diagnose itself.

\subsection{The partition lag and reaction rates}

\begin{definition}[Partition lag]
\label{def:lag}
The \emph{partition lag} $\taup>0$ of an elementary reaction is the
irreducible time during which the reaction's categorical outcome is
undetermined---the minimum time to complete one partition operation. It is
bounded below by the energy--time uncertainty relation,
$\taup\ge\hbar/(2\Delta E)$, and in practice set by the interaction
dynamics (electron transfer $\sim10^{-14}\,$s, enzymatic catalysis
$\sim10^{-6}$--$10^{-3}\,$s, gene expression $\sim10^{2}$--$10^{4}\,$s).
\end{definition}

\begin{theorem}[Reaction rate from partition geometry]
\label{thm:rate}
The rate of a reaction that must cross a partition-depth difference
$\Delta\Depth$ to its transition state is
\[
k\;=\;\frac{1}{\taup}\,b^{-\Delta\Depth}.
\]
\end{theorem}

\begin{proof}
The attempt frequency is $1/\taup$: partition operations occur once per
lag. The success probability is $b^{-\Delta\Depth}$: it is the fraction of
attempts that achieve the required depth reorganisation, since a state of
depth difference $\Delta\Depth$ is one of $b^{\Delta\Depth}$ equiprobable
depth configurations by \Cref{prop:depth-entropy}. The rate is the product.
\end{proof}

\begin{corollary}[Arrhenius and Eyring recovered]
\label{cor:arrhenius}
Writing $E_a=\kB T_{\mathrm{p}}\ln b\cdot\Delta\Depth$ for a characteristic
partition temperature $T_{\mathrm{p}}$, \Cref{thm:rate} is the Arrhenius
form $k=A\,e^{-E_a/RT}$ with pre-exponential $A=1/\taup$, and matches the
Eyring transition-state expression with the activation free energy
identified as the partition-reorganisation cost
\citep{arrhenius1889reaktionsgeschwindigkeit,eyring1935activated}.
\end{corollary}

\begin{proof}
Substitute $E_a=\kB T_{\mathrm{p}}\ln b\cdot\Delta\Depth$ into
$b^{-\Delta\Depth}=e^{-\Delta\Depth\ln b}=e^{-E_a/(\kB T_{\mathrm{p}})}$ and
identify $A=1/\taup$; the Eyring identification of $E_a$ with an activation
free energy follows from \Cref{prop:depth-entropy}
\citep{eyring1935activated,atkins2018physical}.
\end{proof}

\begin{definition}[Catalytic distance]
\label{def:catalytic-distance}
The \emph{catalytic distance} $\dC$ of a reaction is the number of
partition boundaries its mechanism crosses (its categorical distance from
substrate to product state, \Cref{def:catdist}). An enzyme reduces $\dC$ by
providing intermediate partition stages.
\end{definition}

\begin{proposition}[Catalytic efficiency from catalytic distance]
\label{prop:efficiency}
The catalytic efficiency of an enzyme scales inversely with catalytic
distance, $k_{\mathrm{cat}}/K_{\mathrm{M}}\propto 1/(\dC\,\taup)$; reactions
with $\dC=1$ approach the diffusion limit. Equivalently
$\log_{10}(k_{\mathrm{cat}}/K_{\mathrm{M}})\approx c_0-\dC$ for a constant
$c_0$ set by the medium.
\end{proposition}

\begin{proof}
The maximum rate is bounded by substrate encounter (diffusion) and by how
fast one partition operation completes ($1/\taup$). Each additional
boundary ($\dC>1$) requires an additional partition operation, reducing the
rate proportionally; hence efficiency $\propto1/(\dC\,\taup)$
\citep{michaelis1913kinetik,atkins2018physical}.
\end{proof}

\begin{remark}[Monitoring reading]
\label{rem:catalytic-monitoring}
\Cref{prop:efficiency} gives the monitor a handle: an enzyme's observed
catalytic efficiency reports its catalytic distance, and a shift in that
efficiency during a run signals a change in the reaction's partition
mechanism---a mislocalised, inhibited, or misfolded catalyst---before any
bulk concentration has visibly changed.
\end{remark}

\subsection{The reaction network as an information circuit}

The culture's metabolism is a network of reactions; we represent it as a
directed weighted graph and equip it with circuit variables. This is the
structure through which partial measurements will be propagated
(\Cref{sec:completion}).

\begin{definition}[Reaction network]
\label{def:network}
A \emph{reaction network} is a directed weighted graph $G=(\V,\E,w)$ whose
nodes $\V=\{v_1,\dots,v_N\}$ are molecular species and whose directed edges
$e_{ij}$ are elementary reactions, with $w(e_{ij})=k_{ij}$ the forward rate
constant. Each node carries a concentration $[C_i](t)\ge0$. The
\emph{stoichiometric matrix} $\mathbf{S}\in\Z^{N\times R}$ has $S_{ir}$ the
net stoichiometry of species $i$ in reaction $r$, and the dynamics are
$\dot{\mathbf{x}}=\mathbf{S}\,\mathbf{v}(\mathbf{x})$ for reaction-rate
vector $\mathbf{v}$
\citep{heinrich1996regulation,klipp2016systems}.
\end{definition}

\begin{figure*}[!htbp]
    \centering
    \includegraphics[width=\textwidth]{figures/panel2_rates.png}
    \caption{\textbf{Reaction Rates from Partition Geometry Recover Classical
    Kinetics.}
    (\textbf{A})~Partition-geometry rate $k=(1/\taup)\,b^{-\Delta\Depth}$ on a
    log $y$-axis versus the Arrhenius prediction
    $k=A\,e^{-E_a/(\kB T_{\mathrm p})}$ on a log $x$-axis, over $400$ random
    parameter draws (steel-blue), with the identity $y=x$ marked (coral
    dashed). Every point lies on the diagonal: the two expressions are
    identical, not merely correlated, to a maximum relative error below
    $10^{-12}$ over $500$ draws---the exact content of \Cref{cor:arrhenius},
    with $A=1/\taup$ and $E_a=\kB T_{\mathrm p}\ln b\cdot\Delta\Depth$.
    (\textbf{B})~Reaction rate $k$ on a log $y$-axis
    ($\sim\!1$ to $10^{9}\,\mathrm{s}^{-1}$) versus partition-depth
    difference $\Delta\Depth$ on the linear $x$-axis ($0$ to $30$),
    sea-green. The rate falls exactly log-linearly, $k=(1/\taup)b^{-\Delta
    \Depth}$ (\Cref{thm:rate}): each additional partition boundary the
    reaction must cross multiplies the rate by $b^{-1}$.
    (\textbf{C})~Catalytic efficiency $\log_{10}(k_{\mathrm{cat}}/K_{\mathrm
    M})$ on the linear $y$-axis versus catalytic distance $\dC$ on the linear
    $x$-axis ($1$ to $4$), eight enzymes (orange markers) with the
    least-squares fit (blue line). The fitted slope is $-1.31$ with intercept
    $10.06$: efficiency falls roughly one order of magnitude per unit
    catalytic distance, confirming the inverse dependence of
    \Cref{prop:efficiency}; the $d_C=1$ enzymes cluster near the diffusion
    limit and $d_C=4$ near $10^{4}\,\mathrm{M}^{-1}\mathrm{s}^{-1}$.
    (\textbf{D})~Three-dimensional surface of $\log_{10}k$ over the partition
    lag $\log_{10}\taup$ ($-14$ to $-3$) and depth difference $\Delta\Depth$
    ($0$ to $20$), plasma colourmap, viewing angle elevation $24^{\circ}$
    azimuth $-52^{\circ}$. The surface spans $\log_{10}k$ from $-3.0$ to
    $14.0$ and is planar in $(\log_{10}\taup,\Delta\Depth)$: rate is set
    jointly by the attempt frequency $1/\taup$ and the depth barrier, with no
    interaction term, the geometric form of \Cref{thm:rate}.}
    \label{fig:panel2}
\end{figure*}

\begin{definition}[Categorical depth of a species]
\label{def:node-depth}
The \emph{categorical depth} of species $i$ in state $\sigma_i$ is its
self-information $\Hcat_i=-\log_2 P(\sigma_i)$, where $P(\sigma_i)$ is the
stationary probability of that concentration state under the network
dynamics \citep{shannon1948mathematical,cover2006elements}.
\end{definition}

\begin{theorem}[Node potential equals categorical depth]
\label{thm:mu-depth}
For an ideal solution at equilibrium, the chemical potential of species
$i$ satisfies
\[
\mu_i=\mu_i^\circ+RT\ln[C_i]=-\kB T\ln 2\cdot\Hcat_i+c_i,
\]
so the circuit node potential $\phi_i:=\mu_i$ equals the categorical depth
up to the thermal scale $\kB T\ln2$ and a species-specific constant $c_i$.
\end{theorem}

\begin{proof}
The Boltzmann distribution gives
$P(\sigma_i)=Z_i^{-1}\exp(-\mu_i/RT)$ for partition function $Z_i$. Taking
$-\log_2$: $\Hcat_i=\mu_i/(RT\ln2)+\log_2 Z_i$. Solving for $\mu_i$ yields
$\mu_i=RT\ln2\cdot\Hcat_i-RT\ln2\cdot\log_2 Z_i$, which is
$-\kB T\ln2\cdot\Hcat_i+c_i$ once signs and constants are collected against
$\mu_i=\mu_i^\circ+RT\ln[C_i]$ from statistical mechanics
\citep{jaynes1957information,atkins2018physical}.
\end{proof}

\begin{proposition}[Conservation and consistency laws]
\label{prop:kirchhoff}
Assigning to edge $e_{ij}$ the net flux $J_{ij}=k_{ij}[C_i]-k_{ji}[C_j]$,
the network obeys, at steady state, a current law (mass balance) at every
node,
\[
\sum_{j:(j,i)\in\E}J_{ji}-\sum_{j:(i,j)\in\E}J_{ij}=0,
\]
and, around every directed cycle, a voltage law (thermodynamic cycle
consistency),
\[
\sum_{(i,j)\in\mathcal{C}}(\phi_j-\phi_i)=0.
\]
\end{proposition}

\begin{proof}
The current law is $d[C_i]/dt=0$ expanded through
$\sum_r S_{ir}v_r=0$ and regrouped by edge. The voltage law telescopes:
$\sum_{(i,j)\in\mathcal{C}}(\phi_j-\phi_i)=0$ around a closed cycle, and the
concentration-dependent part is $RT\ln\prod_{(i,j)}[C_j]/[C_i]=RT\ln1=0$;
the standard-state part vanishes by thermodynamic consistency of the cycle
\citep{heinrich1996regulation,klipp2016systems}.
\end{proof}

These two laws---mass balance at nodes, thermodynamic consistency around
cycles---are the constraints the monitor propagates measured quantities
through. They hold whatever the sensors report; they are properties of the
network, not of the measurement.

\subsection{Epistemic blindness: why a culture cannot monitor itself}

We close this section with the structural fact that makes external
monitoring necessary. It will reappear, decisively, in the
monitor--control separation of \Cref{sec:separation}.

\begin{definition}[Epistemic blindness]
\label{def:blindness}
A dynamical system $(\V,\mathbf{v})$ is \emph{epistemically blind} to a
partition $\{A,B\}$ of its state space if its flow $\mathbf{v}(\mathbf{x})$
carries no function that correctly classifies $\mathbf{x}$ as belonging to
$A$ or $B$: there is no $f$ with $f(\mathbf{v}(\mathbf{x}))$ deciding
membership for all $\mathbf{x}$.
\end{definition}

\begin{theorem}[A culture is blind to its own state class]
\label{thm:blindness}
Under generic kinetics (mass-action, Michaelis--Menten), a culture running
only forward reaction dynamics is epistemically blind to the partition
\{nominal, aberrant\} of its own states: no function of its instantaneous
flux reliably classifies the culture as nominal or aberrant.
\end{theorem}

\begin{proof}
The rate law $v_{ij}([C_i],[C_j],k_{ij})$ at each edge depends only on
current concentrations and rate constants; it contains no reference to a
stored nominal state. At an aberrant state $\mathbf{x}_D$, the flux
$\mathbf{v}(\mathbf{x}_D)$ encodes only the direction of the next
transition, not the distance $\norm{\mathbf{x}_D-\mathbf{x}_H}$ to a nominal
state $\mathbf{x}_H$---the more so when the aberration lies in altered
parameters (regulation, modification) rather than in concentrations
directly, so that the current flux is indistinguishable from that of a
nominal culture elsewhere in the nominal basin. Hence no function of
$\mathbf{v}(\mathbf{x}_D)$ alone separates aberrant from nominal
\citep{gillespie1977exact,vankampen2007stochastic}.
\end{proof}

\begin{remark}[The necessity of an external observer]
\label{rem:external}
\Cref{thm:blindness} is the reason a bioreactor needs a monitor at all. The
culture propagates whatever state it is in without classifying it; a
faithful \emph{simulation} of the culture inherits exactly this blindness
and so adds no diagnostic information. Detecting that a culture has departed
from its nominal regime requires an observer \emph{outside} the culture's
own dynamics, holding the network topology and a nominal reference against
which the observed state is individuated. The monitoring framework of this
paper \emph{is} that external observer. This is also the first of the three
reasons (\Cref{sec:separation}) that monitoring must be kept separate from
control: a system that acts on the culture is inside the loop it would have
to judge.
\end{remark}

% =====================================================================
\section{Trajectory Completion: Full State from Partial Observation}
\label{sec:completion}
% =====================================================================

A bioreactor is never fully instrumented: a handful of probes report a
handful of the culture's many species, and the rest are unmeasured. The
monitoring task is therefore an \emph{inverse} problem---reconstruct the
complete metabolic state from partial observation---rather than a forward
simulation. We solve it by an iterative procedure that propagates the
measured quantities through the network's own constraints, and we prove it
converges to a unique answer. The state is reported as a bounded region
(honouring the floor of \Cref{sec:floor}), represented as a fuzzy set.

\subsection{Fuzzy state representation}

Because no quantity is known to a point (\Cref{cor:no-point}), we represent
each node's state as a fuzzy membership function over concentration space,
which is exactly a bounded region with graded membership
\citep{zadeh1965fuzzy}.

\begin{definition}[Fuzzy node state]
\label{def:fuzzy}
A \emph{fuzzy state} of node $v_i$ is a function
$\tilde\mu_i\colon\R_{\ge0}\to[0,1]$, where $\tilde\mu_i(c)$ is the degree
to which concentration $c$ is consistent with the available information;
$\tilde\mu_i(c)=1$ means fully consistent, $\tilde\mu_i(c)=0$ means
excluded. Its \emph{support} $\{c:\tilde\mu_i(c)>0\}$ and \emph{core}
$\{c:\tilde\mu_i(c)=1\}$ are its region of possibility and its region of
full consistency. By \Cref{cor:no-point}, the support has width
$\ge\floorc$.
\end{definition}

\begin{definition}[Fuzzy network state; metric]
\label{def:fuzzy-net}
The \emph{fuzzy network state} is the tuple
$\widetilde{\mathbf{X}}=(\tilde\mu_1,\dots,\tilde\mu_N)$. On upper
semi-continuous normal fuzzy sets, the Hausdorff distance between the
$\alpha$-level sets, taken as a supremum over $\alpha\in(0,1]$,
\[
d_H(\tilde\mu,\tilde\nu)=\sup_{\alpha\in(0,1]}
d_{\mathrm{Haus}}\!\bigl([\tilde\mu]^\alpha,[\tilde\nu]^\alpha\bigr),
\]
is a complete metric, and the product metric on the network state is
$d(\widetilde{\mathbf{X}},\widetilde{\mathbf{Y}})=\max_i
d_H(\tilde\mu_i,\tilde\nu_i)$.
\end{definition}

\subsection{Three constraint operators}

The reconstruction is driven by three operators, each tightening the fuzzy
state using a different physical constraint. All three are properties of
the network (\Cref{prop:kirchhoff}) and hold regardless of what the sensors
report.

\begin{definition}[Constraint operators]
\label{def:operators}
On the space of fuzzy network states, define:
\begin{enumerate}[label=(\roman*),leftmargin=2em,nosep]
\item $\mathcal{T}_{\mathrm{bal}}$ (\emph{mass balance}): at each node,
intersect $\tilde\mu_i$ with the set of concentrations consistent with the
node current law (\Cref{prop:kirchhoff}) given the neighbours' current
fuzzy states, propagated by interval arithmetic.
\item $\mathcal{T}_{\mathrm{therm}}$ (\emph{thermodynamic consistency}):
for each independent cycle, restrict node memberships so the fuzzy cycle
potential contains zero in its core (the voltage law).
\item $\mathcal{T}_{\mathrm{kin}}$ (\emph{kinetic feasibility}): for each
node, compute the most-probable causal history of its current state by a
maximum-a-posteriori (Viterbi) trace back through the network
\citep{viterbi1967error}, and restrict $\tilde\mu_i$ to concentrations
whose inferred history is consistent with a nominal reference.
\end{enumerate}
The combined operator is
$\mathcal{T}=\mathcal{T}_{\mathrm{kin}}\circ\mathcal{T}_{\mathrm{therm}}
\circ\mathcal{T}_{\mathrm{bal}}$.
\end{definition}

The kinetic operator uses a fact worth isolating: the most-probable
backward trajectory of a node is a fixed, time-independent property of its
state and the network.

\begin{proposition}[Backward trajectories are time-invariant]
\label{prop:timeinv}
For an autonomous (time-homogeneous) reaction network, the
maximum-a-posteriori causal history of a node's observed state depends only
on that state and the network topology, not on the absolute time of
observation.
\end{proposition}

\begin{proof}
The transition propensities of a mass-action network depend only on the
current molecule counts, not on $t$; the chain is time-homogeneous. Hence
the conditional law of the past trajectory given the present state is
invariant under time-shift, and its maximiser---the MAP history---is the
same whenever the observation is made
\citep{gillespie1977exact,vankampen2007stochastic,viterbi1967error}.
\end{proof}

\begin{remark}[The categorical address]
\label{rem:address}
\Cref{prop:timeinv} lets the monitor treat a node's backward trajectory as
a fixed \emph{address}---a stable identity of its state in the network,
independent of when the culture is sampled. Two nodes at the same
concentration but different network positions have different addresses, so
the monitor distinguishes them even when a bulk sensor cannot. A culture is
nominal when every node's address lies in the nominal basin; an aberration
is a node whose address has escaped it. This is the diagnostic that
\Cref{thm:blindness} showed the culture cannot compute for itself.
\end{remark}

\subsection{Convergence}

\begin{theorem}[Convergence of trajectory completion]
\label{thm:convergence}
The operator $\mathcal{T}=\mathcal{T}_{\mathrm{kin}}\circ
\mathcal{T}_{\mathrm{therm}}\circ\mathcal{T}_{\mathrm{bal}}$ is a
contraction on the complete metric space of fuzzy network states
(\Cref{def:fuzzy-net}). Consequently it has a unique fixed point
$\widetilde{\mathbf{X}}^\ast$, and iteration from any initial state
converges to it geometrically.
\end{theorem}

\begin{proof}
We show each sub-operator is a (weak) contraction and their composition a
strict one. For $\mathcal{T}_{\mathrm{bal}}$, the balance constraint
expresses $[C_i]$ as a convex combination of its in-neighbours' fuzzy
states with coefficients $k_{ji}/\sum_l k_{li}$ that are positive and sum
to at most one; interval propagation and intersection with the prior can
only reduce the Hausdorff distance, so
$d_H(\mathcal{T}_{\mathrm{bal}}(\widetilde{\mathbf X})_i,
\mathcal{T}_{\mathrm{bal}}(\widetilde{\mathbf Y})_i)
\le\sum_{j}\frac{k_{ji}}{\sum_l k_{li}}
d_H(\tilde\mu_j,\tilde\nu_j)$, a contraction with constant $c_1\le1$,
strict once any node is anchored by a measurement. For
$\mathcal{T}_{\mathrm{therm}}$, restricting to the thermodynamic constraint
manifold narrows each interval toward it, giving a factor $c_2<1$
depending on cycle length. For $\mathcal{T}_{\mathrm{kin}}$, restricting to
addresses consistent with the reference reduces support width
monotonically, giving $c_3<1$ whenever the address constraint is
non-trivial. The composition has constant $c\le c_1 c_2 c_3<1$. The space
is complete as a finite product of complete Hausdorff metric spaces
\citep{zadeh1965fuzzy}, so the Banach fixed-point theorem
\citep{banach1922operations} gives a unique fixed point and geometric
convergence $c^n d(\widetilde{\mathbf X}_0,\widetilde{\mathbf X}^\ast)$.
\end{proof}

\begin{corollary}[Sufficient observation yields a bounded region]
\label{cor:observability}
If the measured nodes form a dominating set of the network---every
unmeasured node has a measured neighbour---then the fixed point
$\widetilde{\mathbf X}^\ast$ has every node's support narrowed to a region
of width bounded below by the floor $\floorc$ and above by a width set by
the observation set. The monitor reports this region, not a point.
\end{corollary}

\begin{proof}
Under domination, $\mathcal{T}_{\mathrm{bal}}$ constrains each unmeasured
node from a measured neighbour; $\mathcal{T}_{\mathrm{therm}}$ and
$\mathcal{T}_{\mathrm{kin}}$ remove thermodynamically and kinetically
inconsistent candidates. The fixed-point widths therefore contract to the
region admitted jointly by the three constraints, bounded below by
$\floorc$ (\Cref{thm:floor}) and above by the residual freedom the
observation set leaves.
\end{proof}

\begin{remark}[What the monitor delivers here]
\label{rem:deliver-completion}
\Cref{thm:convergence,cor:observability} are the core monitoring
guarantee for the bulk state: from whatever subset of species the reactor
happens to instrument, the framework reconstructs a bounded, uncertainty-
honest region for \emph{every} species in the network, converging in a
number of iterations logarithmic in the desired precision. Complexity is
$O(N\cdot R\cdot L)$ per iteration for a network of $N$ species, $R$
reactions, and backward-trajectory depth $L$ (\Cref{sec:pipeline}), which
is tractable for genome-scale networks. This is the framework's answer to
the soft-sensor problem, but with a convergence proof and an explicit,
theory-set uncertainty region in place of an empirically fitted error bar.
\end{remark}

\begin{figure*}[!htbp]
    \centering
    \includegraphics[width=\textwidth]{figures/panel3_completion.png}
    \caption{\textbf{Trajectory Completion Contracts to a Unique Bounded
    Region.}
    (\textbf{A})~Iteration residual
    $d(\widetilde{\mathbf X},\widetilde{\mathbf X}_{\mathrm{prev}})$ on a log
    $y$-axis ($10^{-13}$ to $10^{1}$) versus iteration on the linear $x$-axis
    ($0$ to ${\sim}35$), eight random forty-species networks (steel-blue).
    Every curve is a straight line on the semi-log axis: convergence is
    geometric, the signature of a contraction (\Cref{thm:convergence}), each
    network reaching machine precision within $\sim\!30$ passes.
    (\textbf{B})~Distribution of the measured contraction constant $c$ on the
    linear $x$-axis ($0$ to $1.05$) with the boundary $c=1$ marked (coral
    dashed), sea-green. Every measured constant is strictly below one (mean
    $0.50$, maximum $0.63$ over $40$ instances), so the Banach fixed-point
    theorem applies and the fixed point is unique and reached from any
    initial state.
    (\textbf{C})~Distribution of reconstructed interval width normalised by
    the floor, $w/\floorc$, on the linear $x$-axis ($0.9$ to ${\sim}9$),
    orange, with the floor $w/\floorc=1$ marked (coral dashed). A pile-up
    sits exactly at the floor---well-constrained species reach width
    $\floorc$---while distant unmeasured species retain wider intervals; no
    width falls below $\floorc$, the bounded-region guarantee of
    \Cref{cor:observability}. The monitor reports a region, never a point.
    (\textbf{D})~Three-dimensional midpoint trajectories of three tracked
    species from two well-separated initial fuzzy states (init~1 blue,
    init~2 purple) over the coordinates $(x_0,x_1,x_2)$, viewing angle
    elevation $22^{\circ}$ azimuth $-60^{\circ}$; both trajectories converge
    to the identical fixed point (red marker). Distinct starting states reach
    one limit, the uniqueness clause of \Cref{thm:convergence}.}
    \label{fig:panel3}
\end{figure*}

% =====================================================================
\section{Multimodal Reaction Localisation}
\label{sec:localisation}
% =====================================================================

Trajectory completion (\Cref{sec:completion}) reconstructs the bulk state:
how much of each species is present. We now address the finer question of
\emph{where} and \emph{when} an individual reaction occurs. Conventional
bioprocess thinking treats intracellular reaction location as lost to
thermal averaging. We show it is recoverable, because a single reaction
event radiates disturbances into several physically distinct propagation
channels at once, and the intersection of their arrival patterns pins the
event.

\subsection{One event, several modalities}

\begin{definition}[Reaction event; modalities]
\label{def:event}
A \emph{reaction event} occurs at position $\mathbf{r}_0$ and time $t_0$. It
excites, simultaneously, disturbances in several \emph{modalities}, each
governed by distinct propagation physics:
\begin{itemize}[leftmargin=1.4em,nosep]
\item \textbf{chemical}: product species spreading diffusively,
$D\sim10^{-11}\,\mathrm{m}^2/\mathrm{s}$;
\item \textbf{acoustic}: a pressure wave from molecular rearrangement,
$c_s\sim1540\,\mathrm{m/s}$;
\item \textbf{thermal}: reaction-enthalpy heat spreading diffusively,
$\alpha\sim10^{-7}\,\mathrm{m}^2/\mathrm{s}$;
\item \textbf{electromagnetic}: charge redistribution, screened over the
Debye length $\lambda_D\sim0.5$--$1\,\mathrm{nm}$;
\item \textbf{vibrational}: a local bond-frequency change, extent
$\sim0.1\,\mathrm{nm}$;
\item \textbf{categorical}: a discrete change in partition coordinates,
counted exactly (\Cref{def:catdist}), with no threshold.
\end{itemize}
\end{definition}

Each modality reaches an observation point $\mathbf{r}$ at a characteristic
time, defining an arrival-time function
$\mathcal{T}_i(\mathbf{r};\mathbf{r}_0,t_0)$ whose level sets are
\emph{arrival-time surfaces}. Ballistic modalities (acoustic) give surfaces
expanding linearly, $|\mathbf{r}-\mathbf{r}_0|=c_s(t-t_0)$; diffusive
modalities (chemical, thermal) give surfaces expanding as $\sqrt{t}$,
$|\mathbf{r}-\mathbf{r}_0|=2\sqrt{D\mathcal{L}(t-t_0)}$ for a logarithmic
threshold factor $\mathcal{L}$; screened and local modalities (EM,
vibrational) give hard constraint spheres; the categorical modality gives a
discrete, threshold-free constraint
\citep{crank1979mathematics,smoluchowski1917versuch,herzberg1950spectra,%
boxer2009stark}.

\subsection{The intersection principle}

\begin{theorem}[Localisation by intersection]
\label{thm:intersection}
Given accurate arrival times of at least three modalities with distinct
propagation laws, observed at a sufficient set of points, the reaction
coordinates $(\mathbf{r}_0,t_0)$ are uniquely determined. The probability
of a false localisation decreases exponentially in the number of
modalities.
\end{theorem}

\begin{proof}
Consider the acoustic (linear-in-distance) and chemical (quadratic-in-
distance) constraints observed at $\mathbf{r}$. Writing
$r=|\mathbf{r}-\mathbf{r}_0|$, the acoustic constraint gives
$r=c_s(t_A-t_0)$ and the chemical constraint gives
$r^2=4D\mathcal{L}(t_C-t_0)$. Eliminating $r$ yields a quadratic in $t_0$
whose physical (earlier) root is selected by causality, fixing $t_0$ and
hence $r$. A third distinct-law modality, or observations at further
points, resolves direction by trilateration
\citep{fang2008simple}. Because the modalities obey functionally
independent laws (linear vs.\ quadratic vs.\ hard-sphere vs.\ discrete),
their constraint surfaces meet transversally, so coincident false
intersections form a measure-zero set; with independent measurement noise
across $N$ modalities the false-localisation probability falls as
$e^{-\Theta(N)}$.
\end{proof}

\begin{theorem}[Resolution multiplies with modalities]
\label{thm:resolution}
If modality $i$ excludes a fraction $\eps_i\in(0,1)$ of the candidate
volume, then the combined localisation resolution is
\[
\delta r_{\mathrm{comb}}=\delta r_{\mathrm{single}}\prod_{i=1}^{N}
\eps_i^{1/3},
\]
the $1/3$ exponent arising from three-dimensional intersection.
\end{theorem}

\begin{proof}
For statistically independent modalities the admissible volume is
$V_0\prod_i\eps_i$; the characteristic linear dimension scales as
$V^{1/3}$, giving $\delta r_{\mathrm{comb}}=\delta
r_{\mathrm{single}}(\prod_i\eps_i)^{1/3}$
\citep{fang2008simple}.
\end{proof}

\begin{example}[Order of magnitude]
\label{ex:resolution}
With six modalities of typical exclusion factors $\eps_i$ between
$10^{-2}$ and $10^{-6}$, the product $\prod_i\eps_i^{1/3}\sim10^{-7}$, so a
diffraction-limited single-channel resolution of $\sim300\,\mathrm{nm}$
sharpens toward the sub-nanometre scale, in practice limited to
$\sim0.1$--$1\,\mathrm{nm}$ by thermal noise. The gain is the payoff for
fusing physically independent channels rather than refining one.
\end{example}

\subsection{Opacity-independent measurement}

The categorical modality plays a special role, because it lives in the
space of distinctions, not the space of physical positions.

\begin{proposition}[Categorical distance is independent of physical
distance and opacity]
\label{prop:opacity}
The categorical distance $\dcat$ (\Cref{def:catdist}) between two states is
mathematically independent of their spatial separation and of the optical
opacity of the intervening medium.
\end{proposition}

\begin{proof}
$\dcat$ is defined on partition coordinates, which reside in the space of
distinctions; spatial distance and optical opacity are defined on physical
coordinates. The two coordinate systems are constructed independently, so
their generated distances are uncorrelated: a state may be categorically
close yet physically distant or occluded, and vice versa.
\end{proof}

\begin{figure*}[!htbp]
    \centering
    \includegraphics[width=\textwidth]{figures/panel4_localisation.png}
    \caption{\textbf{Multimodal Intersection Localises a Reaction; Resolution
    Multiplies.}
    (\textbf{A})~Localisation error on a log $y$-axis (nm) versus the number
    of intersecting modalities on the linear $x$-axis ($1$ to $3$), median
    with inter-quartile bars over $120$ events per point (steel-blue). Error
    falls monotonically as modalities are added---from ${\sim}1.6\,\mu$m with
    the acoustic channel alone to $\sim\!20\,$nm with acoustic, thermal, and
    chemical combined---the intersection principle of \Cref{thm:intersection}:
    functionally independent arrival-time surfaces meet at a single event.
    (\textbf{B})~Combined resolution on a log $y$-axis (nm) versus modality
    count on the linear $x$-axis ($1$ to $6$), sea-green, with the
    $1\,$nm line marked (coral dashed). The product-law factor
    $\prod_i\eps_i^{1/3}$ drives resolution from $\sim\!300\,$nm toward the
    sub-nanometre scale, confirming \Cref{thm:resolution}: each independent
    channel contributes a multiplicative $\eps_i^{1/3}$ to the resolving
    power.
    (\textbf{C})~Median localisation error on the linear $y$-axis (nm) versus
    optical opacity of the intervening medium on the linear $x$-axis
    ($0$ to $50$), orange. The error is flat---${\sim}20\,$nm at every
    opacity---because the timing and categorical channels do not read
    photons: localisation is invariant to opacity, the content of
    \Cref{prop:opacity} and \Cref{cor:opacity}, so a reaction behind an
    opaque barrier remains monitorable.
    (\textbf{D})~Three-dimensional scatter of true (blue) and recovered (red
    crosses) event positions over the $10\,\mu$m domain
    $(x,y,z)$, viewing angle elevation $20^{\circ}$ azimuth $-58^{\circ}$.
    Recovered points coincide with the true positions to sub-nanometre in the
    noiseless limit (maximum error ${\sim}10^{-16}\,$m), so the two markers
    overlie one another: the intersection is exact
    (\Cref{thm:intersection}).}
    \label{fig:panel4}
\end{figure*}

\begin{corollary}[Subsurface reactions remain measurable]
\label{cor:opacity}
A reaction occurring behind an optically opaque barrier---inside a cell,
behind a membrane, within an aggregate---remains categorically accessible
to the monitor, provided its partition signature is distinguishable. The
optically inaccessible reaction is not thereby monitor-inaccessible.
\end{corollary}

\begin{proof}
Immediate from \Cref{prop:opacity}: opacity governs photon transmission,
$\dcat$ governs distinguishability, and the two are independent.
\end{proof}

\begin{remark}[The reaction pathway as a chromatographic process]
\label{rem:chromatography}
There is a productive way to picture the localisation channels
concretely. Imagine the culture---or a sampled stream from it---traversing
a structured medium (a packed column of conductive dielectric beads) under
flow, optical access, and a small electrical bias. Each reacting species,
as it passes, modulates three physical channels at once: it changes the
local viscous resistance (mechanical), the local refractive index
(optical), and the local conductance across bead junctions (electrical).
The species is then operationally a \emph{partition operator}: each
channel-completing or channel-interrupting event is a categorical
transition, and the temporal sequence of such events is the reaction being
read out. The result is a multi-channel fingerprint of the reaction
whose resolving power multiplies across channels
(\Cref{thm:resolution})---the same intersection principle, realised as an
instrument in which separation and detection are one event rather than two.
A standard mass spectrometer is the degenerate special case that reads only
one channel. This reading is offered as orientation and as a route to a
physical instrument; the theorems above stand on the abstract modalities
alone.
\end{remark}

% =====================================================================
\section{Distributed Monitoring by a Phase-Locked Swarm}
\label{sec:swarm}
% =====================================================================

A single monitoring agent identifies the state of a local region. To
monitor a whole vessel---and to identify the culture as completely as
possible---one deploys many agents. The question is how they coordinate. We
show that agents which synchronise to a common cellular oscillation
coordinate at zero communication overhead once coherent, that their
resolving power \emph{increases} with agent count, and that the loss of
their coherence is itself a first-class fault signal.

\subsection{Agents as coupled oscillators}

\begin{definition}[Monitoring agent]
\label{def:agent}
A \emph{monitoring agent} is characterised by a natural oscillation
frequency $\omega_i=2\pi f_i$, a phase $\phi_i(t)$, and a
\emph{cell registry} mapping each cell of a fixed partition of the
observable space to a pre-compiled action. Agents communicate only by
emitting bare timing pulses: the sole communicated quantity is \emph{when}
a pulse arrives, never a semantic payload. A physical perturbation shifts
an agent's inter-pulse interval into a different cell, triggering the
associated action, with no explicit encoding or decoding.
\end{definition}

\begin{definition}[Coupling; order parameter]
\label{def:kuramoto}
Agents evolve by the Kuramoto dynamics
\citep{kuramoto1984chemical,strogatz2000kuramoto}
\[
\dot\phi_i=\omega_i+\frac{K}{n}\sum_{j=1}^{n}\sin(\phi_j-\phi_i),
\]
with coupling strength $K\ge0$. The natural frequencies are drawn from a
symmetric unimodal distribution of standard deviation $\sigma_\omega$. The
collective coherence is the \emph{order parameter}
\[
\Rens\,e^{i\psi}=\frac1n\sum_{j=1}^{n}e^{i\phi_j},\qquad \Rens\in[0,1],
\]
with $\Rens=0$ full incoherence and $\Rens=1$ perfect synchrony.
\end{definition}

\begin{theorem}[Critical coupling]
\label{thm:critical}
There is a critical coupling $K_c=2\sigma_\omega/\pi$ such that for $K<K_c$
the only stable state is incoherent ($\Rens^\ast=0$), while for $K>K_c$ a
coherent state $\Rens^\ast=\sqrt{1-K_c/K}$ emerges through a supercritical
pitchfork bifurcation.
\end{theorem}

\begin{proof}
In the mean-field limit each agent obeys
$\dot\phi_i=\omega_i+K\Rens\sin(\psi-\phi_i)$; splitting into locked
($|\omega_i|\le K\Rens^\ast$) and drifting populations and imposing
self-consistency on $\Rens^\ast$ gives, near the bifurcation,
$1=\tfrac{\pi}{2}K^2 g(0)$ for the frequency density $g$; substituting
$g(0)$ for a unimodal distribution of spread $\sigma_\omega$ yields
$K_c=2\sigma_\omega/\pi$, and the normal-form reduction gives the stable
branch $\Rens^\ast=\sqrt{1-K_c/K}$
\citep{strogatz2000kuramoto,pikovsky2001synchronization}.
\end{proof}

\subsection{Zero-overhead coordination}

\begin{theorem}[Zero-overhead coordination]
\label{thm:zero-overhead}
In the phase-locked regime $\Rens\ge0.95$, with a cell partition whose
minimum cell width exceeds the residual phase spread
$2\pi(1-\Rens)/\bar\omega$, the per-iteration coordination overhead is
identically zero: every agent assigns the same cell to each event and hence
computes the same action, with no message exchange.
\end{theorem}

\begin{proof}
At $\Rens\ge0.95$ all agents rotate at a common angular velocity up to a
residual phase error $|\delta_i|\le\arccos(\Rens)$; the induced spread of
inter-pulse intervals is bounded by $\eps_0=2\pi(1-\Rens)/\bar\omega$. If
the minimum cell width exceeds $\eps_0$, any two agents' intervals for the
same event fall in the same cell, so their (identical, pre-distributed)
registries yield the same action. No exchange of cell assignments is needed
because the assignments already agree; hence the coordination overhead is
zero \citep{lynch1996distributed}.
\end{proof}

\begin{remark}[Contrast with content-based consensus]
\label{rem:consensus}
\Cref{thm:zero-overhead} is the coordination analogue of a coherent
collective ground state: below the decoherence threshold, agents act as one
without reconciliation events, whereas content-based consensus protocols
require $O(\log n)$ to $O(n^2)$ messages per decision
\citep{lynch1996distributed}. The saving is structural, not algorithmic:
there is nothing to reconcile because there is no semantic content in the
channel, only timing.
\end{remark}

\subsection{More agents mean finer identification}

\begin{theorem}[Composition inflates distinguishability]
\label{thm:composition}
A federation whose agents contribute $D$ independent sensing channels in
total can distinguish, over $m$ timing events,
\[
T(m,D)=D\,(D+1)^{m-1}
\]
joint states. Adding an agent with $d$ channels multiplies the count by
$\bigl((D+d+1)/(D+1)\bigr)^{m-1}>1$.
\end{theorem}

\begin{proof}
Model the joint state at each event as a draw from an alphabet of size
$D+1$ ($D$ channel-excited states plus one nominal). The first event has
$D$ informative states; each later event draws independently from all
$D+1$, giving $D(D+1)^{m-1}$. Replacing $D$ by $D+d$ gives the multiplier,
which exceeds one for $d\ge1$
\citep{cover2006elements}.
\end{proof}

\begin{remark}[The identification goal]
\label{rem:identify}
\Cref{thm:composition} is why adding monitoring agents \emph{helps} rather
than merely costs: distinguishability grows super-exponentially in the
number of channels, reversing the usual curse of dimensionality. Combined
with the resolution multiplication of \Cref{thm:resolution} and the
never-completable medium of \Cref{def:medium}, this makes precise the
monitoring goal: identify the culture as completely as the deployed
channels allow, knowing the floor forbids ever completing it. More agents
resolve more of the culture; none of them finish it.
\end{remark}

\subsection{Decoherence as a fault signal}

\begin{theorem}[Graceful decoherence and detection bound]
\label{thm:decoherence}
If the effective coupling falls below $K_c$ (through channel degradation or
a genuine process disturbance), the coherence $\Rens$ relaxes toward zero at
rate $(K_c-K)/2$, and the departure is detectable within
$m^\ast=\lceil1/\eps\rceil$ timing events for timing precision $\eps$. The
federation does not fail abruptly; it degrades to the next-lower
coordination regime, continuing to monitor at correspondingly higher
overhead.
\end{theorem}

\begin{proof}
For $K<K_c$ the coherent fixed point disappears and $\Rens$ decays with
rate set by the curvature of the coupling potential at the origin,
$(K_c-K)/2$; a per-event decrease of at least $\eps$ is observable after at
most $1/\eps$ events. The five coordination regimes are ordered by
$\Rens$, so crossing a boundary triggers the fallback protocol of the next
regime rather than collapse
\citep{strogatz2000kuramoto,pikovsky2001synchronization}.
\end{proof}

\begin{corollary}[Swarm decoherence is a monitoring output]
\label{cor:decoherence-output}
The loss of swarm coherence is itself a first-class, early fault
signal---detectable within $\lceil1/\eps\rceil$ events, before any
individual agent has diagnosed a local aberration---and is reported by the
monitor as such.
\end{corollary}

\begin{proof}
By \Cref{thm:decoherence} the coherence responds to any disturbance that
perturbs the agents' common clock, and the response is detectable within a
bounded number of events; the monitor exposes $\Rens$ and its regime as an
output.
\end{proof}

\begin{remark}[Security by absence of content]
\label{rem:security}
Because the inter-agent channel carries only timing and no semantic
payload, there is no content to inject, corrupt, or forge: a spurious pulse
can at worst shift an interval to an adjacent cell, selecting a different
but still pre-approved action, and a replayed pulse is rejected by a
monotone phase counter. The monitoring swarm is thus structurally immune to
content-injection attacks, a property inherited from the timing-only
channel rather than from any cryptographic layer.
\end{remark}


% =====================================================================
\section{Monitoring Coincidence Nodes: States True of the Population,
Present in No Cell}
\label{sec:coincidence}
% =====================================================================

We now reach a capability that distinguishes this framework from every
per-cell or bulk-average monitoring method: the ability to monitor a
process that occurs only under a \emph{conjunction} of conditions that no
single cell need ever satisfy at once. Such a process is invisible to
conventional monitoring---there is no cell to probe while it happens, and it
is averaged away in any bulk signal---yet it is individuable from the
population's joint evidence. The information about it is true of the culture
as a structured whole while being present in no member of it.

\subsection{Coincidence nodes}

\begin{definition}[Condition; coincidence node]
\label{def:coincidence}
A \emph{condition} is a predicate on the observable state---for example
``$\mathrm{pH}=x$,'' ``$T=y$,'' ``the UV--vis spectrum lies in band $z$.''
A \emph{coincidence node} is a network node (a reaction or state,
\Cref{def:network}) whose occurrence requires the \emph{conjunction} of two
or more conditions $c_1\wedge c_2\wedge\cdots\wedge c_p$, each supplied by a
distinct evidentiary channel. Its \emph{physical support} at a given instant
is the set of cells that satisfy all $p$ conditions simultaneously; the
support may be empty even when each condition is satisfied by many cells
individually.
\end{definition}

\begin{example}[A triple-conjunction node]
\label{ex:triple}
Let a node $\nu$ fire only when a cell is simultaneously at pH $x$,
temperature $y$, and exhibiting UV--vis spectrum $z$. In a well-mixed
reactor, many cells are at pH $x$, many at temperature $y$, and many showing
spectrum $z$, but the fraction at all three at once may be vanishingly small
or momentarily zero. No probe placed on any cell reliably catches $\nu$; no
bulk average contains it, since the overwhelming majority of cells
contributing to any average are not at the intersection. By the standards of
conventional monitoring, $\nu$ is unmonitorable.
\end{example}

\subsection{Individuation of an empty-support node}

The framework individuates $\nu$ regardless, because individuation is by
negation against the medium (a property of the whole,
\Cref{thm:negation})---not by finding a cell that instantiates $\nu$.

\begin{definition}[Categorical address of a coincidence node]
\label{def:coincidence-address}
The \emph{categorical address} of a coincidence node $\nu=c_1\wedge\cdots
\wedge c_p$ is the point of the state space at which all $p$ conditions'
constraint regions intersect, expressed in the categorical coordinates of
\Cref{sec:sspace}. It is defined by the conditions, not by any cell: it is
the categorical intersection of the $p$ conditions, whether or not a cell
occupies it.
\end{definition}

\begin{theorem}[Conjunction individuation]
\label{thm:conjunction}
Let $\nu=c_1\wedge\cdots\wedge c_p$ be a coincidence node, each condition
$c_i$ witnessed by an independent channel that individuates a constraint
region of separation cost $\ge\floorc$ (\Cref{thm:floor}). If the
categorical address of $\nu$ (\Cref{def:coincidence-address}) is reachable
from the committed evidence---that is, the $p$ constraint regions have a
non-empty categorical intersection---then $\nu$ is individuable, and the
monitor reports its reconstructed state-region, \emph{even when the physical
support of $\nu$ is empty}.
\end{theorem}

\begin{proof}
Each condition $c_i$ is individuated against the medium by negation
(\Cref{thm:negation}), fixing a constraint region $\Reac_i$ of the state
space of separation cost $\ge\floorc$ (\Cref{thm:floor}); this fixing uses
only the channel witnessing $c_i$ and the medium, and in particular does not
require a cell that satisfies $c_i$ to also satisfy the others. The address
of $\nu$ is the categorical intersection $\bigcap_{i=1}^{p}\Reac_i$
(\Cref{def:coincidence-address}). Because categorical distance is
independent of physical distance and of whether any cell occupies the
intersection (\Cref{prop:opacity}), the reachability of the address is a
statement about the intersection of the $\Reac_i$ in categorical
coordinates, not about the physical support of $\nu$. When that intersection
is non-empty, the trajectory-completion operators of \Cref{sec:completion}
propagate the $p$ constraints---each committed by a different channel---to a
unique fixed-point region for $\nu$ (\Cref{thm:convergence}), of width
$\ge\floorc$ (\Cref{cor:observability}). This region is $\nu$'s
reconstructed state. No step required a cell instantiating all $p$
conditions; hence $\nu$ is individuated with empty physical support.
\end{proof}

\begin{figure*}[!htbp]
    \centering
    \includegraphics[width=\textwidth]{figures/panel5_coincidence.png}
    \caption{\textbf{A Coincidence Node: True of the Population, Present in No
    Cell.}
    (\textbf{A})~Number of cells satisfying each single condition and the
    three-way conjunction, on a symmetric-log $y$-axis, over a population of
    $5{,}000$ cells (bars; the conjunction bar in coral). Each condition
    holds for many cells---$3{,}361$ at the pH band, $3{,}291$ at the
    temperature band, $3{,}348$ at the spectral band---yet the conjunction is
    satisfied by \emph{zero} cells: the physical support is empty, so the
    node is unmonitorable by any per-cell method (\Cref{thm:unmonitorable}).
    (\textbf{B})~The reconstructed categorical address as normalised
    condition ranges for the three axes (sea-green horizontal bars, each
    axis rescaled to $[0,1]$ of its physiological range). The three
    condition regions have a non-empty, bounded intersection---the address
    of the node---even though no cell occupies it, the object individuated by
    \Cref{thm:conjunction} and \Cref{def:coincidence-address}.
    (\textbf{C})~Node detected ($0$ or $1$) by three methods: per-cell,
    bulk-average, and conjunction individuation (bars; failing methods in
    coral, the succeeding method in green). Only conjunction individuation
    returns the node; the per-cell and bulk routes cannot, the separation
    proved in \Cref{thm:unmonitorable}.
    (\textbf{D})~Three-dimensional scatter of the population in condition
    space (pH $5.5$ to $8.5$, temperature $30$ to $42$, spectral band $0$ to
    $1$), each cell coloured by how many of the three bands it satisfies
    ($0$ grey, $1$ blue, $2$ orange---never three), viewing angle elevation
    $20^{\circ}$ azimuth $-62^{\circ}$. The cloud forms a cross whose centre
    is hollow; the coral wireframe box marks the empty categorical address at
    the band intersection---the state the monitor reconstructs from the
    population's joint evidence, present in no member of it
    (\Cref{cor:populational}).}
    \label{fig:panel5}
\end{figure*}

\begin{corollary}[The information is populational, not cellular]
\label{cor:populational}
The state of a coincidence node is carried by the joint structure of the
population's evidence---one condition per channel, drawn across different
cells and different instants---and by no single cell. It is information that
is true for the culture as a whole and present in no member of it.
\end{corollary}

\begin{proof}
By \Cref{thm:conjunction} the reconstruction of $\nu$ uses one committed
constraint per condition, each witnessed by a distinct channel and,
generically, sourced from a different subpopulation of cells; no cell
supplies more than a proper subset of the conjunction (else the support
would be non-empty at that cell). The individuated region for $\nu$ is
therefore a function of the joint evidence across channels and cells, and is
recoverable from no single cell's state.
\end{proof}

\subsection{Why conventional monitoring cannot do this}

\begin{theorem}[Coincidence nodes are unmonitorable per cell or in bulk]
\label{thm:unmonitorable}
A coincidence node with empty physical support cannot be detected by any
method that reads individual cells or bulk population averages, yet is
individuable by conjunction individuation (\Cref{thm:conjunction}).
\end{theorem}

\begin{proof}
A per-cell method detects a node only by observing a cell that instantiates
it; with empty physical support there is no such cell, so per-cell detection
fails. A bulk-average method reports a population mean; the contribution of a
node instantiated by a measure-zero (or empty) fraction of cells is
dominated by the overwhelming complementary fraction not at the
intersection, so the node is not resolvable from the average
(\Cref{thm:blindness} gives the sharper statement that the bulk flux carries
no classifier of the node's occurrence). Both conventional routes therefore
fail. Conjunction individuation, by contrast, reconstructs the node from the
categorical intersection of independently committed conditions
(\Cref{thm:conjunction}), which does not require the node's physical support
to be non-empty. Hence the node is monitorable by the framework and by
neither conventional route.
\end{proof}

\begin{remark}[Monitoring processes that occur after specific sequences]
\label{rem:sequence}
The conjunction of \Cref{def:coincidence} need not be simultaneous in the
naive sense: a condition may itself be ``the culture has passed through
state $s$,'' encoded as a committed step in the monotone history of the
completion procedure. A coincidence node may then require an \emph{ordered}
sequence of conditions---a process that occurs only after a specific
succession of events. Because the backward addresses of
\Cref{prop:timeinv} record committed history, the framework individuates
such sequence-gated nodes by the same theorem: the conditions are the
history-predicates, and their categorical intersection is the address of the
sequence-triggered process. Processes that ``only occur after this, then
this, then this'' are thus monitorable exactly as spatial conjunctions are.
\end{remark}

\begin{remark}[Latent nodes and the identification goal]
\label{rem:latent}
\Cref{thm:conjunction} instructs the monitor to carry coincidence nodes as
\emph{latent but individuable}: a coincidence node is represented as a
categorical address with a reconstructed region, flagged with its current
physical support (possibly empty), whether or not any cell currently
occupies it. This is the sharpest form of the identification goal
(\Cref{rem:identify}): the monitor accumulates not only what each cell is
doing, but what is true of the population's joint structure---coincidences
and sequences that no cell exhibits alone. It is knowledge extracted from
the culture as a whole, and it is exactly the knowledge that conventional
monitoring, bound to cells and averages, cannot reach.
\end{remark}


% =====================================================================
\section{Closure: When Monitoring Is Entitled to Report}
\label{sec:closure}
% =====================================================================

A monitor that fuses many channels must decide when it has seen enough. We
give a stopping rule---\emph{closure}---that is strictly stronger than
crossing a confidence threshold, prove it terminates, and show it yields
exactly two honest outcomes: a single reported state, or an explicit
declaration that the evidence is contested.

\subsection{Evidence, catalysts, and coherence}

\begin{definition}[Catalyst]
\label{def:catalyst}
A \emph{catalyst} is an evidence-producing operation---a probe reading, a
spectroscopic inversion, a completed trajectory reconstruction, a
localisation---that moves the monitor's estimate toward the state being
identified by committing a contact in the contact graph. Catalysts of any
physical source (in-line probe, off-line assay, inferred quantity) are
\emph{uniform}: the framework encodes only the cost of the evidence they
produce, never its physical origin (\Cref{def:contact}).
\end{definition}

\begin{definition}[Catalytic power; composition]
\label{def:power}
The \emph{power} $\kappa\in[0,1]$ of a catalyst is the fraction of the
above-floor uncertainty it removes. A chain of catalysts of powers
$\kappa_1,\dots,\kappa_n$ has composite power
$\kappa=1-\prod_{i=1}^{n}(1-\kappa_i)$.
\end{definition}

\begin{proposition}[Diminishing returns and diversification]
\label{prop:diminishing}
Repeating one catalyst of power $\kappa$ gives composite power
$1-(1-\kappa)^n$, strictly increasing toward but never reaching $1$; a
chain drives the residual uncertainty to zero only if
$\sum_i\kappa_i=\infty$. Hence independent catalysts addressing distinct
evidentiary channels saturate faster and more robustly than repeated near-
duplicate ones.
\end{proposition}

\begin{proof}
The residual fraction after $n$ catalysts is $\prod_{i\le n}(1-\kappa_i)$;
for a single repeated catalyst this is $(1-\kappa)^n\to0$ but never $0$ at
finite $n$. Taking logarithms, $\sum_i\log(1-\kappa_i)\to-\infty$ iff
$\sum_i\kappa_i=\infty$ (comparison test, since $-\log(1-\kappa)\asymp
\kappa$ for small $\kappa$). Distinct channels contribute independent
powers and so accumulate the divergent sum faster than duplicates, whose
marginal power collapses \citep{cover2006elements}.
\end{proof}

\begin{theorem}[Coherence requires three independent catalysts]
\label{thm:triangle}
A monitored state is robustly grounded---its identification survives the
failure of any single catalyst---if and only if its supporting catalysts
contain a strongly connected cycle of at least three mutually reinforcing,
independently sourced members. One or two catalysts do not suffice: a
single source has no internal check, and two sources can only mutually
assert, with no third to break a disagreement.
\end{theorem}

\begin{proof}
\emph{Necessity.} An acyclic support structure has a terminal catalyst with
no incoming support; removing it, or any catalyst supporting only
downstream members, collapses the chain below it, so no acyclic structure
is robust to single removal. A one-cycle is self-support (vacuous); a
two-cycle is mutual assertion between exactly two members, and removing
either leaves one unsupported, failing a majority check. The shortest
robust support cycle therefore has length $\ge3$.
\emph{Sufficiency.} In a strongly connected triangle each member is
supported by the other two; removing one leaves the remaining two still
mutually supporting, so the identification persists.
\end{proof}

\begin{remark}[Monitoring reading]
\label{rem:triangle-monitoring}
\Cref{thm:triangle} is the framework's formalisation of a familiar
laboratory instinct: a critical process reading should be corroborated by
at least three independent lines of evidence---say an in-line probe, a
spectroscopic inference, and a completed mass balance---before it is
trusted, precisely because a lone reading (or two mutually confirming ones)
cannot survive the failure of one source. The requirement is structural and
decidable from the \emph{signs} of pairwise support alone, without knowing
the powers.
\end{remark}

\subsection{The closure criterion}

\begin{definition}[Reachable state set]
\label{def:reachable}
At a given monitoring stage, let $\mathcal{C}$ be the set of state-regions
reachable by some admissible chain of catalysts invoked so far, taken up to
the equivalence ``indistinguishable by any test of the reported state.''
\end{definition}

\begin{definition}[Closure]
\label{def:closure}
A monitoring query is \emph{closed} if, for every catalyst available but
not yet invoked, extending the query by that catalyst cannot add a new
region to $\mathcal{C}$: every further reachable outcome already lies in a
region present in $\mathcal{C}$.
\end{definition}

\begin{theorem}[Closure is strictly stronger than a confidence threshold]
\label{thm:closure-stronger}
There exist monitoring configurations and thresholds $\theta$ in which a
reading meets the confidence criterion (its estimate exceeds $\theta$)
while the query is not closed: a further, not-yet-invoked catalyst would
reach a region in a distinct equivalence class.
\end{theorem}

\begin{proof}
Take two disjoint, internally self-consistent clusters $A$ and $B$ of
supporting evidence, each reachable by a single catalyst $\gamma_A$,
$\gamma_B$. A chain via $\gamma_A$ alone terminates at a region in $A$ and
satisfies any fixed threshold trivially, since it is internally consistent.
Yet $\gamma_B$, not yet invoked, reaches a region in $B$ distinct from that
in $A$. The confidence criterion is met after $\gamma_A$; the query is not
closed, because $\gamma_B$ can still add a new region.
\end{proof}

\begin{remark}[Why closure is correct]
\label{rem:closure-correct}
\Cref{thm:closure-stronger} isolates the standard failure of threshold
stopping: one early, internally consistent sensor reports high confidence
while a second, independent sensor would conclude otherwise. Closure
forbids stopping until \emph{every} available further catalyst has been
checked to land in an already-reached region. It is not that the current
reading looks confident; it is that the space of readings has stopped
producing new destinations. For a bioreactor this is the difference between
``the DO probe is confident'' and ``no available measurement could change
our conclusion.''
\end{remark}

\subsection{Two honest outcomes}

\begin{theorem}[Convergent closure or honest decline]
\label{thm:decline}
Every monitoring query over a finite catalyst set terminates, in finitely
many steps, in exactly one of two states:
\begin{enumerate}[label=(\roman*),leftmargin=2.2em,nosep]
\item \textbf{convergent closure}: $\mathcal{C}$ stabilises to a single
equivalence class, and the monitor reports that region as the identified
state;
\item \textbf{contested closure (decline)}: $\mathcal{C}$ stabilises but
contains more than one class, and the monitor reports \emph{decline}---no
single sufficient state exists---together with the distinct regions found.
\end{enumerate}
\end{theorem}

\begin{proof}
The catalyst set is finite, so the reachable state set is finite and can
grow only finitely often before every catalyst has been tried against every
reached class; at that point \Cref{def:closure} holds by exhaustion, so
closure is reached in finitely many steps. If the stabilised $\mathcal{C}$
has one class, outcome (i) holds; if more than one, no single region is
licensed and outcome (ii) holds. The two cases are exhaustive and mutually
exclusive on a non-empty partition.
\end{proof}

\begin{remark}[Decline is information, not failure]
\label{rem:decline}
\Cref{thm:decline}(ii) is essential to an honest monitor. When independent
lines of evidence disagree---an in-line probe says one thing, an off-gas
inference another---the correct output is neither to loop forever nor to
silently emit one of the contested readings as if it were unique, but to
report that the sources disagree, with the alternatives. Contested closure
localises exactly where independent inquiry diverges, which is itself
valuable diagnostic output and, frequently, the earliest sign of a
developing fault.
\end{remark}


% =====================================================================
\section{The Separation of Monitoring from Control}
\label{sec:separation}
% =====================================================================

We now establish, as a structural result rather than a design preference,
that process monitoring (the framework of this paper) must be kept as a
\emph{separate} framework from process control. Three independent arguments
force the separation, and together they fix the worst-case behaviour of the
monitor as its most valuable property.

\subsection{Three arguments for separation}

\begin{theorem}[Monitoring requires an observer outside the loop]
\label{thm:sep-blindness}
A system that acts on the culture cannot serve as an unbiased observer of
the culture it acts on.
\end{theorem}

\begin{proof}
By \Cref{thm:blindness}, diagnosing a culture's state class requires an
observer external to the culture's own forward dynamics, holding the network
topology and a nominal reference. A control element that actuates the
culture is, by construction, part of the loop whose state is in question:
its actions enter the very dynamics the diagnosis must judge, and its
observations are of a culture it has already perturbed. Fusing observation
with actuation therefore returns the observer to the inside of the loop and
reintroduces the blindness \Cref{thm:blindness} identifies. Hence unbiased
monitoring requires the observer to be separated from the actuator.
\end{proof}

\begin{theorem}[The floor forbids the point-precision control presumes]
\label{thm:sep-floor}
A controller acting on a monitored quantity as if it were known to a point
acts on precision the framework forbids.
\end{theorem}

\begin{proof}
By \Cref{cor:no-point}, every monitored quantity resolves to a region of
width $\ge\floorc>0$, never a point. A control law that consumes a point
estimate discards the region and the residual, acting on a precision that
\Cref{thm:floor} shows cannot exist. A correct architecture hands the
controller the honest region---floor, residue, and, where evidence is
contested, the decline of \Cref{thm:decline}(ii)---which is a different
interface from a point estimate and belongs to a distinct framework whose
job is to decide what to do given a region, not a number.
\end{proof}

\begin{theorem}[Observation terminates at knowing, not acting]
\label{thm:sep-obs}
The monitoring operation is closed under identification and does not extend
to actuation.
\end{theorem}

\begin{proof}
Every operation in the framework---individuation (\Cref{thm:negation}),
completion (\Cref{thm:convergence}), localisation (\Cref{thm:intersection}),
closure (\Cref{thm:decline})---terminates in a reported state-region or a
decline. None produces an action on the culture: the output type of the
framework is information (regions, addresses, coherence, fault signals),
not actuation. Actuation is a distinct operation the framework never
performs, and joining it would require an output type---a command---that no
theorem of the framework produces.
\end{proof}

\subsection{The worst case is a safe state}

\begin{theorem}[Worst-case monitoring output]
\label{thm:worst-case}
If the control framework is absent, misconfigured, or declines to act, the
monitoring framework still delivers, for the culture: a bounded state-region
per species (\Cref{cor:observability}), localised reaction events
(\Cref{thm:intersection}), a swarm-coherence health signal
(\Cref{cor:decoherence-output}), and, per query, either an identified region
or an explicit decline (\Cref{thm:decline}). It degrades to detailed,
uncertainty-bounded information about the process and nothing about it
breaks.
\end{theorem}

\begin{proof}
Each cited result is a property of the monitoring framework alone and
invokes no control action: completion, localisation, coherence estimation,
and closure all run on observation only. Removing or disabling control
removes no input to any of them. Hence the monitor's outputs persist
unchanged, and its worst case is the delivery of information with no action
taken.
\end{proof}

\begin{corollary}[Independent failure modes]
\label{cor:independent-failure}
Under the separation, a failure of control does not induce a failure of
monitoring: the monitor continues to produce information, leaving the system
in the safe state of \Cref{thm:worst-case}. A fused monitor--controller has
no such guarantee, since a fault in its control logic disables its
observation with it.
\end{corollary}

\begin{proof}
By \Cref{thm:worst-case} the monitor's outputs are independent of control;
by construction a fused system shares a failure surface between observation
and actuation, so a control fault removes observation. The separated
architecture therefore strictly dominates the fused one in failure
containment.
\end{proof}

\begin{remark}[The design principle]
\label{rem:design-principle}
The three theorems and their corollary give a single design principle: the
deliverables of this framework end at information---bounded regions,
categorical addresses, coherence and fault signals, and decline reports---
handed \emph{across} an interface to a separate control framework whose task
is to decide actions given that information and its uncertainty. The
separation is not caution; it is what the finiteness of the observer, the
resolution floor, and the closure of observation under identification
jointly require. Its dividend is that the monitor's worst case is never a
wrong action, only detailed knowledge of the process---which is exactly the
floor one wants a monitoring system to guarantee.
\end{remark}

\begin{figure*}[!htbp]
    \centering
    \includegraphics[width=\textwidth]{figures/panel6_swarm.png}
    \caption{\textbf{Swarm Synchronisation, Zero-Overhead Locking, and
    Decoherence as a Fault Signal.}
    (\textbf{A})~Kuramoto order parameter $\Rens$ on the linear $y$-axis
    ($0$ to $1$) versus coupling ratio $K/K_c$ on the linear $x$-axis
    ($0.3$ to $8$), measured (blue markers) against the mean-field branch
    $\Rens^\ast=\sqrt{1-K_c/K}$ (coral line), with the critical point
    $K/K_c=1$ marked (green dotted). Coherence stays near zero below
    $K_c=2\sigma_\omega/\pi\approx0.637$ and rises above it, converging to
    the mean-field prediction at strong coupling (absolute error $0.04$ at
    $8K_c$), the transition of \Cref{thm:critical}; the near-critical gap is
    the expected finite-size lag.
    (\textbf{B})~Distinguishable joint-state count $T(m,D)=D(D+1)^{m-1}$ on a
    log $y$-axis versus channel count $D$ on the linear $x$-axis
    ($1$ to $40$), orange, at $m=10$ events. The count rises
    super-exponentially---$2.4\times10^{10}$ at $D=10$,
    $1.1\times10^{20}$ at $D=100$---confirming \Cref{thm:composition}:
    adding channels multiplies distinguishability, reversing the usual curse
    of dimensionality.
    (\textbf{C})~Order parameter $\Rens$ on the linear $y$-axis ($0$ to $1$)
    versus time on the linear $x$-axis, steel-blue, with the coupling drop
    marked (coral dashed). Held above $K_c$ the swarm locks at
    $\Rens\approx0.9$; when the coupling falls below $K_c$ the coherence
    relaxes to ${\sim}0.05$, the graceful-decoherence fault signal of
    \Cref{thm:decoherence} and \Cref{cor:decoherence-output}.
    (\textbf{D})~Three-dimensional surface of the order parameter over
    coupling $K/K_c$ ($0.5$ to $6$) and network size $\log_{10}n$
    ($1.7$ to $2.9$), viridis colourmap, viewing angle elevation
    $24^{\circ}$ azimuth $-56^{\circ}$. The synchronisation manifold is the
    ridge where $\Rens$ climbs from near-zero to near-one across $K_c$; the
    transition sharpens with $n$, approaching the thermodynamic-limit
    bifurcation of \Cref{thm:critical}.}
    \label{fig:panel6}
\end{figure*}


% =====================================================================
\section{The Complete Monitoring Pipeline}
\label{sec:pipeline}
% =====================================================================

We assemble the results into a single procedure and state its correctness
and complexity. The pipeline takes the reactor's sensor stream and produces,
at each monitoring cycle, an uncertainty-bounded report of the culture's
state together with any localised events and fault signals---and never an
action.

\subsection{The procedure}

\begin{algorithm}[t]
\caption{Bioreactor Monitoring Cycle}
\label{alg:monitor}
\begin{algorithmic}[1]
\Require Reaction network $G=(\V,\E,w)$ with standard potentials; partial
sensor readings $\mathcal{O}$; nominal reference addresses
$\{\mathcal{A}_i^H\}$; catalyst registry; swarm phase stream
$\{\phi_j\}$; floor $\floorc$.
\Ensure State report $\widetilde{\mathbf X}^\ast$ (bounded region per
species); localised events; coherence $\Rens$ and regime; per-query
verdict (identified region or \textsc{decline}).
\State \textbf{Coordinates.} Map each sensor reading to S-entropy
coordinates $\Scoord=(\Sk,\St,\Se)$ (\Cref{def:sspace}); individuate each
against the medium by negation (\Cref{thm:negation}).
\State \textbf{Completion.} Initialise fuzzy states: measured nodes as
narrow regions of width $\ge\floorc$, unmeasured nodes as maximum-entropy
priors over physiological ranges.
\Repeat
   \State $\widetilde{\mathbf X}\gets\mathcal{T}_{\mathrm{bal}}
          (\widetilde{\mathbf X})$
          \Comment{mass balance at every node}
   \State $\widetilde{\mathbf X}\gets\mathcal{T}_{\mathrm{therm}}
          (\widetilde{\mathbf X})$
          \Comment{thermodynamic cycle consistency}
   \State $\widetilde{\mathbf X}\gets\mathcal{T}_{\mathrm{kin}}
          (\widetilde{\mathbf X})$
          \Comment{Viterbi backward-address feasibility}
\Until{$d(\widetilde{\mathbf X},\widetilde{\mathbf X}_{\mathrm{prev}})
       <\eps$}
       \Comment{contraction, \Cref{thm:convergence}}
\State \textbf{Health.} For each node, compare its backward address to the
nominal $\mathcal{A}_i^H$; flag \emph{address escape} where the categorical
distance exceeds the per-node threshold (\Cref{rem:address}).
\State \textbf{Localisation.} For each event of interest, intersect the
arrival-time surfaces of the available modalities to recover
$(\mathbf r_0,t_0)$ (\Cref{thm:intersection}); resolution multiplies with
modality count (\Cref{thm:resolution}).
\State \textbf{Coherence.} Estimate the swarm order parameter $\Rens$ from
$\{\phi_j\}$; classify the coordination regime; if $\Rens$ has fallen a
regime, raise the decoherence fault (\Cref{cor:decoherence-output}).
\State \textbf{Closure.} For each monitoring query, invoke catalysts until
no available catalyst adds a new state-region (\Cref{def:closure}); require
a $\ge3$-catalyst coherence cycle for a grounded verdict
(\Cref{thm:triangle}).
\If{$\mathcal{C}$ is a single class}
   \State \Return identified region (convergent closure)
\Else
   \State \Return \textsc{decline} with the distinct regions (contested
   closure, \Cref{thm:decline})
\EndIf
\State \textbf{(No actuation.)} Hand the report across the interface to the
separate control framework (\Cref{sec:separation}).
\end{algorithmic}
\end{algorithm}

\subsection{Correctness and complexity}

\begin{theorem}[Pipeline correctness]
\label{thm:pipeline-correct}
Each monitoring cycle of \Cref{alg:monitor} terminates and returns a report
in which (i) every species is assigned a bounded region of width
$\ge\floorc$; (ii) every localised event is unique with false-localisation
probability $e^{-\Theta(N_{\mathrm{mod}})}$; (iii) every query ends in a
single identified region or an explicit decline; and (iv) no output is an
action on the culture.
\end{theorem}

\begin{proof}
The completion loop terminates and is correct by \Cref{thm:convergence}
(contraction to a unique fixed point) and \Cref{cor:observability} (bounded
regions of width $\ge\floorc$), giving (i). Localisation uniqueness and the
false-localisation bound are \Cref{thm:intersection}, giving (ii). Closure
termination in one of two outcomes is \Cref{thm:decline}, giving (iii). No
step produces an action: every return value is a region, a signal, or a
decline, by \Cref{thm:sep-obs}, giving (iv).
\end{proof}

\begin{proposition}[Per-cycle complexity]
\label{prop:complexity}
For a network of $N$ species, $R$ reactions, maximum degree $\Delta$,
backward-trajectory depth $L$, $N_{\mathrm{mod}}$ localisation modalities,
$M$ observation points, and $n$ swarm agents, one monitoring cycle costs
\[
O\!\bigl(N R L\,\log(1/\eps)\bigr)\ (\text{completion})
\;+\;O\!\bigl(M^2 N_{\mathrm{mod}}\bigr)\ (\text{localisation})
\;+\;O(n)\ (\text{coherence}),
\]
with completion dominating for genome-scale networks; the completion
sub-steps parallelise across nodes.
\end{proposition}

\begin{proof}
Completion: mass balance is $O(N\Delta)$ per iteration, cycle consistency
$O(R\Delta)$, and Viterbi addresses $O(N R L)$; geometric convergence
(\Cref{thm:convergence}) needs $O(\log(1/\eps)/\log(1/c))$ iterations,
giving $O(N R L\log(1/\eps))$. Localisation is $O(M^2 N_{\mathrm{mod}})$ for
pairwise arrival-time constraints across modalities and points, plus a
bounded-iteration least-squares refinement. Coherence is $O(n)$ for the
order-parameter sum. Node updates in completion are independent given fixed
neighbours, so they parallelise
\citep{viterbi1967error,fang2008simple}.
\end{proof}

\begin{remark}[Reading the pipeline]
\label{rem:pipeline}
\Cref{alg:monitor} is the framework in operational form. Every step is a
theorem discharged: coordinates and negation
(\Cref{sec:sspace,sec:observer}), completion (\Cref{sec:completion}),
addresses and health (\Cref{rem:address}), localisation
(\Cref{sec:localisation}), coherence (\Cref{sec:swarm}), closure
(\Cref{sec:closure}), and the terminal absence of actuation
(\Cref{sec:separation}). What the operator receives each cycle is a
complete, uncertainty-honest picture of the culture---bulk state, event
locations, and health---with an explicit statement of what remains
contested, and nothing that acts on the vessel.
\end{remark}


% =====================================================================
\section{Mapping the Framework to a Physical Bioreactor}
\label{sec:mapping}
% =====================================================================

The framework is stated in abstract quantities; we now identify each with a
concrete bioreactor observable, so that a reader can see how the pipeline
would be instantiated on a real vessel. Nothing in this section adds a
claim; it is a dictionary.

\subsection{Coordinates and the state space}

The knowledge coordinate $\Sk$ is set by the spread of the culture's
identifiable state---operationally, the entropy of the reconstructed
concentration distribution over the network's species. The temporal
coordinate $\St$ is set by the spread of characteristic process timescales,
from fast metabolic turnover to slow expression dynamics. The evolution
coordinate $\Se$ is set by the uncertainty in the culture's trajectory---
whether it is heading toward, holding at, or departing from its nominal
operating regime. Each is a normalised information measure computed from the
sensor stream and the completed state.

\subsection{Sensors as catalysts}

Every physical sensor is a catalyst in the sense of \Cref{def:catalyst},
producing evidence that individuates part of the culture against the medium.
Table~\ref{tab:sensors} gives representative correspondences. The framework
treats them uniformly: what matters is the cost and independence of the
evidence, not the transducer.

\begin{table}[t]
\centering
\caption{Representative physical sensors as catalysts, and the modality
(\Cref{def:event}) each primarily excites.}
\label{tab:sensors}
\small
\begin{tabular}{@{}lll@{}}
\toprule
\textbf{Physical sensor} & \textbf{Evidence produced} & \textbf{Primary modality} \\
\midrule
Dissolved-oxygen probe      & O\textsubscript{2} activity, respiration & chemical \\
pH probe                    & proton activity, acid--base state        & chemical/electrical \\
Off-gas mass spectrometry   & \ce{CO2}/\ce{O2} exchange, respiratory quotient & chemical \\
Dielectric (capacitance)    & viable biomass, membrane state           & electrical \\
Raman / NIR spectroscopy    & metabolite and product bands             & vibrational/optical \\
Calorimetry (heat flux)     & metabolic heat, reaction enthalpy        & thermal \\
Acoustic / ultrasonic       & density, rearrangement signatures        & acoustic \\
Fluorescence (NAD(P)H)      & redox and electronic state               & optical \\
\bottomrule
\end{tabular}
\end{table}

\subsection{The oxygen sensor field}

Molecular oxygen has a distinctive role. It is paramagnetic, ubiquitous in
an aerobic culture ($\sim10^2\,\mu\mathrm{M}$, of order $10^8$ molecules per
cell), highly mobile, and its vibrational frequency shifts with local
electric field through the Stark effect, $\delta\omega\propto|\mathbf E|^2$
\citep{herzberg1950spectra,boxer2009stark}. It therefore constitutes a
dense, distributed field-sensing network already present in the medium: an
enormous number of simultaneous, in-situ reporters of local charge
redistribution. Two consequences follow for the pipeline. First, dissolved
oxygen is not merely a respiration proxy but a spatial field probe, feeding
the electromagnetic and vibrational modalities of localisation
(\Cref{sec:localisation}). Second, a characteristic cellular oscillation
carried on the oxygen field furnishes the common clock to which a
monitoring swarm phase-locks (\Cref{sec:swarm})---the coupling medium is
physically the same signal the culture runs on, so coherence and its loss
report directly on the culture, not on an artificial synchronisation layer.

\subsection{The cell as a three-layer electrostatic system}

For the electromagnetic modality, it is useful to picture the cell
concretely as a three-layer electrostatic system: a fixed negative charge
on the genomic DNA (a sequence-independent scaffold of order $-1\,$nC), a
mobile-ion cytoplasm, and a dynamically charged plasma membrane. This
configuration sustains a static cytoplasmic field of order
$10^5$--$10^6\,\mathrm{V/m}$
\citep{israelachvili2011intermolecular,hille2001ion}, against which reaction
events register as local charge redistributions---the electromagnetic
disturbances of \Cref{def:event}. The picture is offered to make the EM
modality tangible; the localisation theorems do not depend on it.

\subsection{Health readouts}

The framework exposes three complementary health signals, all computed from
observation alone. The \emph{address-escape} signal (\Cref{rem:address})
flags any species whose reconstructed causal history has left the nominal
basin---an early, concentration-independent indicator of regulatory or
mechanistic departure. The \emph{swarm-coherence} signal
(\Cref{cor:decoherence-output}) flags a loss of collective synchronisation
across the monitoring array within $\lceil1/\eps\rceil$ events. The
\emph{contested-closure} signal (\Cref{thm:decline}(ii)) flags, per query,
that independent evidence lines disagree---frequently the earliest sign of
a developing fault, and precisely the situation in which a fused
monitor--controller would be most dangerous.


% =====================================================================
\section{Numerical Validation}
\label{sec:validation}
% =====================================================================

The theorems of this paper are proved from the stated premises and require
no computation. It is nonetheless useful to exhibit them on explicitly
constructed instances, both as a check against error and to display the
quantitative shape of each result. This section reports a validation suite
of seven independent experiments, one per major result, comprising
thirty-four individual checks. Every check passed; the suite is
deterministic given a master seed and was confirmed to pass under
independent re-seeding, so the outcomes are not artefacts of a particular
random draw. \Cref{tab:validation} summarises the suite, and
\Cref{fig:panel1,fig:panel2,fig:panel3,fig:panel4,fig:panel5,fig:panel6,%
fig:panel7} display the measured quantities against the predicted bounds.

\begin{table}[t]
\centering
\caption{Validation suite. Each experiment tests the theorems of one
section by direct numerical experiment; all checks pass.}
\label{tab:validation}
\small
\begin{tabular}{@{}llc@{}}
\toprule
\textbf{Experiment} & \textbf{Result tested} & \textbf{Checks} \\
\midrule
Resolution floor      & \Cref{thm:floor}, \Cref{cor:no-point}, \Cref{prop:refine} & 4/4 \\
Reaction rates        & \Cref{thm:rate}, \Cref{cor:arrhenius}, \Cref{prop:efficiency} & 3/3 \\
Trajectory completion & \Cref{thm:convergence}, \Cref{cor:observability} & 5/5 \\
Localisation          & \Cref{thm:intersection}, \Cref{thm:resolution}, \Cref{prop:opacity} & 5/5 \\
Coincidence nodes     & \Cref{thm:conjunction}, \Cref{thm:unmonitorable} & 7/7 \\
Swarm coordination    & \Cref{thm:critical}, \Cref{thm:zero-overhead}, \Cref{thm:composition} & 4/4 \\
Closure \& coherence  & \Cref{thm:triangle}, \Cref{thm:closure-stronger}, \Cref{thm:decline} & 6/6 \\
\midrule
\textbf{Total} & & \textbf{34/34} \\
\bottomrule
\end{tabular}
\end{table}

\subsection{The resolution floor}

We built random contact graphs (three to fourteen parts, a medium vertex
adjacent to every part, edge weights drawn at or above a floor
$\floorc=1$) and computed each part's separation cost
$\sigma(v)=\min\text{-cut}(v,\med)$ exactly by maximum flow. Across
$\sim\!2{,}600$ parts of $300$ instances, every $\sigma(v)\ge\floorc$
(\Cref{thm:floor}): the minimum observed ratio $\sigma(v)/\floorc$ was
$1.02$, and no separation was free (\Cref{cor:no-point}). Refining a graph
by adding parts and contacts of weight $\ge\floorc$ never produced a
sub-floor separation, and the floor remained strictly positive as the
medium grew from four to sixty-four parts (\Cref{prop:refine}).
\Cref{fig:panel1} displays the result.



\subsection{Reaction rates}

The partition-geometry rate $k=(1/\taup)\,b^{-\Delta\Depth}$ was checked
against the Arrhenius form $k=A\,e^{-E_a/(\kB T_{\mathrm p})}$ with
$A=1/\taup$ and $E_a=\kB T_{\mathrm p}\ln b\cdot\Delta\Depth$ over $500$
random parameter draws: the two agreed to a maximum relative error below
$10^{-12}$ (\Cref{cor:arrhenius}). Against a panel of eight enzymes spanning
catalytic distances $\dC\in\{1,\dots,4\}$, the fitted slope of
$\log_{10}(k_{\mathrm{cat}}/K_{\mathrm M})$ in $\dC$ was $\approx-1$ (one
order of magnitude per partition boundary) with a strong negative rank
correlation, confirming the inverse dependence of \Cref{prop:efficiency}.
\Cref{fig:panel2} displays the result.


\subsection{Trajectory completion}

On random reaction networks of forty species with a dominating measured
set, the combined constraint operator was iterated from arbitrary initial
fuzzy states. The measured contraction constant was strictly below one on
every instance, iteration converged geometrically, and two distinct
initial states reached the same fixed point (\Cref{thm:convergence}); every
reconstructed interval had width $\ge\floorc$ and contained the truth at
measured nodes (\Cref{cor:observability}). \Cref{fig:panel3} displays the
result.



\subsection{Multimodal localisation}

A reaction event was localised by least-squares intersection of acoustic
(linear-in-distance), chemical, and thermal (quadratic-in-distance)
arrival-time surfaces observed at twelve points in a $10\,\mu\mathrm m$
domain. Noiseless recovery was exact to sub-nanometre; under realistic
timing noise the median error stayed sub-micron (\Cref{thm:intersection}).
Adding modalities reduced the error monotonically, and the product-law
factor $\prod_i\eps_i^{1/3}$ reached the sub-nanometre scale
(\Cref{thm:resolution}); the error was invariant to a varied optical
opacity (\Cref{prop:opacity}). \Cref{fig:panel4} displays the result.



\subsection{Coincidence nodes}

A population of $5{,}000$ cells was arranged in a three-dimensional
condition space (pH, temperature, spectral band) so that each condition
held for many cells while \emph{no} cell satisfied all three at once---an
empty physical support. A per-cell detector returned nothing and a
bulk-average detector could not report the truth, yet the categorical
address (the intersection of the three condition regions) was non-empty and
bounded, and conjunction individuation reconstructed the node's region
(\Cref{thm:conjunction}, \Cref{thm:unmonitorable}). The ordered-sequence
variant behaved identically. A control with non-empty support confirmed
that per-cell detection succeeds precisely when the support is non-empty.
\Cref{fig:panel5} displays the result.



\subsection{Swarm coordination}

Simulating the Kuramoto model with Lorentzian natural frequencies, the
order parameter stayed near zero below $K_c=2\sigma_\omega/\pi$ and rose
toward $\sqrt{1-K_c/K}$ above it, matching the mean-field branch at strong
coupling (\Cref{thm:critical}). In the phase-locked regime the residual
phase spread fell below the cell width, so all agents shared one cell with
zero coordination messages (\Cref{thm:zero-overhead}); the distinguishable
joint-state count $T(m,D)=D(D+1)^{m-1}$ grew super-exponentially with
channel count (\Cref{thm:composition}); and dropping the coupling below
$K_c$ relaxed the coherence downward, the decoherence fault signal
(\Cref{thm:decoherence}). \Cref{fig:panel6} displays the result.



\subsection{Closure and coherence}

The composite catalytic power $1-\prod_i(1-\kappa_i)$ matched sequential
residual reduction to machine precision and stayed strictly below one
(\Cref{prop:diminishing}). A support structure was robust to single-catalyst
removal exactly when it contained a mutually-supporting group of at least
three members---never at one or two (\Cref{thm:triangle}). A two-cluster
configuration met a confidence threshold while remaining un-closed
(\Cref{thm:closure-stronger}), and every query terminated in convergent
closure or honest decline (\Cref{thm:decline}). \Cref{fig:panel7} displays
the result.

\begin{figure*}[!htbp]
    \centering
    \includegraphics[width=\textwidth]{figures/panel7_closure.png}
    \caption{\textbf{Catalytic Power, the Coherence Triangle, and Closure
    Beyond Confidence.}
    (\textbf{A})~Composite catalytic power computed by the closed form
    $1-\prod_i(1-\kappa_i)$ on the linear $y$-axis versus sequential residual
    reduction on the linear $x$-axis ($0$ to $1$), over $400$ random catalyst
    chains (steel-blue), with the identity $y=x$ marked (coral dashed). Every
    point lies on the diagonal (maximum absolute error $0.0$): the
    multiplicative composition law of \Cref{prop:diminishing} is exact, and
    the composite power stays strictly below one.
    (\textbf{B})~Robustness to single-catalyst removal ($0$ or $1$) versus
    mutual-support group size on the linear $x$-axis ($1$ to $6$), bars with
    the $2\!\to\!3$ boundary marked (orange dotted); groups of size one and
    two are never robust (coral), groups of three and above always are
    (green). Robust grounding first appears at three mutually reinforcing
    catalysts and never below, the exact content of \Cref{thm:triangle}: a
    lone or paired source cannot survive the failure of one member.
    (\textbf{C})~Confidence value and closure status on the linear $y$-axis
    ($0$ to $1$) versus threshold $\theta$ on the linear $x$-axis
    ($0.5$ to $0.99$) for two-cluster configurations: cluster-A confidence
    (blue) stays at $\sim\!0.95$, above every threshold, while the closure
    flag (coral) stays at zero, with the gap shaded. A configuration can meet
    any confidence threshold yet remain not closed---an un-invoked catalyst
    reaches a distinct region---the strict-separation result of
    \Cref{thm:closure-stronger}.
    (\textbf{D})~Three-dimensional surface of the two-catalyst composite
    power $1-(1-\kappa_1)(1-\kappa_2)$ over the unit square
    $\kappa_1,\kappa_2\in[0,1]$, magma colourmap, viewing angle elevation
    $22^{\circ}$ azimuth $-55^{\circ}$. The surface rises from zero at the
    origin to one at $(1,1)$, saturating fastest along either axis and
    slowest on the diagonal where both catalysts are weak---the geometry of
    diminishing returns and the value of diverse, independent evidence
    (\Cref{prop:diminishing}).}
    \label{fig:panel7}
\end{figure*}

\begin{remark}[Status of the numerical witness]
\label{rem:witness-status}
The figures confirm the results on concrete instances and display their
quantitative form; they are evidence of correctness and magnitude, not a
substitute for the proofs, which hold for every instance satisfying the
premises. Where a result is asymptotic---the mean-field synchronisation
branch of \Cref{thm:critical}---the suite asserts what holds at finite size
(incoherence below $K_c$, coherent onset above it, convergence to the
prediction at strong coupling) and records the expected near-critical
deviation rather than concealing it. Where the reference data are real and
noisy---enzyme efficiencies against catalytic distance---the suite tests the
theorem's actual claim, the inverse trend, not a strict monotonicity the
scatter would not support. All thirty-four checks passed.
\end{remark}


% =====================================================================
\section{Discussion, Scope, and Conclusion}
\label{sec:discussion}
% =====================================================================

\subsection{Relation to existing monitoring practice}

The framework does not discard established bioprocess monitoring; it
supplies the foundation that established practice lacks and recovers its
components as special cases. Flux balance analysis
\citep{orth2010fba} is the mass-balance operator
$\mathcal{T}_{\mathrm{bal}}$ of \Cref{def:operators} without the
thermodynamic and kinetic constraints; the framework adds cycle consistency
and backward-address feasibility and, crucially, an existence-and-
uniqueness proof for the reconstructed state (\Cref{thm:convergence}).
Soft sensors \citep{luttmann2012soft,mandenius2015bioprocess} are
instances of trajectory completion, but with the empirically fitted error
bar replaced by a theory-set uncertainty region bounded below by the
resolution floor (\Cref{sec:floor}). Multivariate statistical monitoring
\citep{wold1987principal,qin2012survey} detects departures in a learned
subspace; the framework detects them as address escape in the network's own
coordinates, with a mechanistic rather than a correlational reading, and
adds the swarm-coherence and contested-closure signals that have no
statistical-monitoring analogue. Classical estimation
\citep{kalman1960new} propagates a point estimate with a covariance; the
framework propagates a bounded region and reports, honestly, when the region
cannot be reduced to one class. The Process Analytical Technology programme
\citep{fda2004pat} calls for exactly the uncertainty-explicit, mechanism-
grounded monitoring the framework provides.

\subsection{Why the framework is not phenomenology}

A reasonable concern is whether the framework merely re-labels known
results. It does not: from the single premise of a finite observer against a
non-completable medium, it \emph{derives} the resolution floor
(\Cref{thm:floor}), the convergence of state reconstruction
(\Cref{thm:convergence}), the localisation and resolution laws
(\Cref{thm:intersection,thm:resolution}), the coordination and coherence
results (\Cref{thm:zero-overhead,thm:decoherence}), the closure criterion
(\Cref{thm:decline}), and the monitor--control separation
(\Cref{thm:sep-blindness,thm:sep-floor,thm:sep-obs}). Each is a consequence
of the premise, not an assumption added to fit a case. Where the framework
recovers classical results---Arrhenius and Eyring kinetics
(\Cref{cor:arrhenius}), Kirchhoff-type conservation laws
(\Cref{prop:kirchhoff})---it does so as corollaries, which is evidence for
the premise rather than a restatement of the corollaries.

\subsection{Scope and limitations}

The framework is a monitoring theory and a monitoring algorithm; it is not
a control law, a completed software system, or an experimental validation on
a physical vessel, and the separation of monitoring from control
(\Cref{sec:separation}) is deliberate. Several assumptions bound its present
reach. The trajectory-completion convergence (\Cref{thm:convergence})
assumes an autonomous, time-homogeneous network; strongly time-varying
kinetics (rapid induction, phase transitions in the culture) require a
non-autonomous extension that would forgo the time-invariance of backward
addresses (\Cref{prop:timeinv}). The localisation theorems
(\Cref{sec:localisation}) assume a homogeneous propagation medium and a
single event at a time; spatial heterogeneity and simultaneous events
require, respectively, spatially varying propagation parameters and source
separation. The swarm results (\Cref{sec:swarm}) assume controllable
coupling and an all-to-all topology in the thermodynamic limit; finite-size
and heterogeneous-topology corrections are needed for small arrays. The
categorical modality's contribution to localisation depends on maintaining
distinguishable partition signatures over the relevant scale. None of these
limitations touches the resolution floor or the monitor--control separation,
which follow from the premise alone.

\subsection{Conclusion}

We have developed, from the single premise that a bioreactor is a population
of finite observers against a non-completable medium, a self-contained
framework for bioprocess monitoring. The framework answers the three
questions with which we began. A bioreactor measurement is the individuation
of a bounded region of the culture's state space, fixed by negation against
the medium; it is never a point. The precision of monitoring is bounded below
by a resolution floor that is a property of the finite culture and not of the
instrument, so an honest monitor reports regions of width no less than that
floor. And a monitoring query is entitled to report when it reaches
closure---when no available further evidence can change its outcome---
terminating either in a single identified region or in an explicit
declaration that the evidence is contested.

Around these answers we built an operational pipeline: reconstruct the full
metabolic state from partial observation by a provably convergent completion
of the network's own constraints; localise individual reactions by
intersecting the arrival-time surfaces of several physically independent
modalities, with resolution that multiplies as modalities are added and that
is indifferent to optical opacity; coordinate a distributed monitoring swarm
at zero overhead through phase-locking to the culture's own oxygen-borne
clock, reading the loss of coherence as an early fault; and terminate every
query at closure, with contested closure exposing exactly where independent
evidence diverges.

Finally, we established that monitoring must be kept separate from control,
because a system that acts on the culture cannot be its unbiased observer,
because the resolution floor forbids the point-precision a fused controller
would presume, and because observation is closed under identification and
does not extend to actuation. The dividend of the separation is the
framework's worst-case guarantee: with control absent or failed, the monitor
still delivers detailed, uncertainty-bounded information about the process
and never a wrong action. In the worst case, one simply has thorough
knowledge of the bioreactor---which is precisely what a monitoring system
should be built to guarantee.


% =====================================================================
\section*{Acknowledgements}
% =====================================================================
The author thanks the Technical University of Munich for institutional
support. All results are proved from the stated premises; material that
could not be established at that standard has been omitted rather than
weakened.

% =====================================================================
\bibliographystyle{plainnat}
\bibliography{references}
% =====================================================================

\end{document}
